% League Simulator -- Handbuch
% Bauen: latexmk -lualatex league-simulator-handbuch.tex   (siehe README.md)
\documentclass[11pt,a4paper,parskip=half,toc=bibliography,numbers=noenddot]{scrreprt}

\usepackage[ngerman]{babel}
\usepackage{fontspec}
\usepackage{unicode-math}
\setmainfont{Source Serif Pro}
\setsansfont{Inter}[Scale=MatchLowercase]
\setmonofont{Source Code Pro}[Scale=MatchLowercase]
\setmathfont{STIX Two Math}[Scale=MatchLowercase]

\usepackage{microtype}
\usepackage{csquotes}
\usepackage{xcolor}
\usepackage{booktabs,tabularx,longtable,array}
\usepackage{graphicx}
\usepackage{amsmath}
\usepackage{mathtools}
\usepackage{tikz}
\usetikzlibrary{arrows.meta,positioning,fit,backgrounds,shapes.geometric,calc}
\usepackage{listings}
\usepackage[most]{tcolorbox}
\usepackage{enumitem}
\usepackage{scrlayer-scrpage}
\usepackage[locale=DE,group-minimum-digits=4,group-separator={\,}]{siunitx}

% Screenshots aus bilder/ (PNG). Fehlt die Datei, steht ein Platzhalter da.
\newcommand{\screenshot}[2]{%
  \begin{figure}[!htbp]\centering
  \IfFileExists{bilder/#1.png}%
    {\fcolorbox{rule}{white}{\includegraphics[width=0.94\textwidth,height=0.5\textheight,keepaspectratio]{bilder/#1.png}}}%
    {\fcolorbox{rule}{surface2}{\parbox[c][4cm][c]{0.9\textwidth}{\centering\sffamily\small\color{ink3}Screenshot \texttt{bilder/#1.png}}}}%
  \caption{#2}\label{fig:#1}
  \end{figure}}
\usepackage{xurl}
\usepackage[hidelinks]{hyperref}
\usepackage[htt]{hyphenat}
\renewcommand{\_}{\BreakableUnderscore}

% --- Farben aus RCode/site_assets/site.css -----------------------------
\definecolor{ink}{HTML}{15130F}
\definecolor{ink2}{HTML}{312E29}
\definecolor{ink3}{HTML}{6B655C}
\definecolor{rule}{HTML}{D4CEBF}
\definecolor{rulesoft}{HTML}{E3DED0}
\definecolor{surface2}{HTML}{F5F3EE}
\definecolor{akzent}{HTML}{B0261E}
\definecolor{tinte}{HTML}{1F3A4D}
\definecolor{ocker}{HTML}{B8893A}
\definecolor{gruen}{HTML}{3F7D3A}

\addtokomafont{disposition}{\color{ink}}
\addtokomafont{chapter}{\rmfamily\mdseries\Huge}
\addtokomafont{section}{\rmfamily\mdseries\Large}
\addtokomafont{subsection}{\rmfamily\mdseries\large}
\addtokomafont{subsubsection}{\rmfamily\bfseries\normalsize}
\addtokomafont{part}{\rmfamily\mdseries}
\addtokomafont{partnumber}{\sffamily\small\color{ink3}}
\addtokomafont{chapterprefix}{\sffamily\small\color{akzent}}
\renewcommand*{\chapterformat}{\mbox{\textcolor{akzent}{\rule[.2em]{2em}{1.5pt}}\enskip\thechapter\enskip}}
\addtokomafont{pageheadfoot}{\sffamily\footnotesize\color{ink3}}
\pagestyle{scrheadings}
\clearpairofpagestyles
\ohead{\leftmark}
\ofoot*{\pagemark}
\automark[chapter]{chapter}

\setlist{itemsep=2pt,topsep=4pt}
\setlength{\emergencystretch}{2em}

% --- Kästen --------------------------------------------------------------
\newtcolorbox{merke}[1][]{enhanced,breakable,colback=surface2,colframe=surface2,
  borderline west={3pt}{0pt}{ocker},arc=0pt,left=10pt,right=8pt,top=6pt,bottom=6pt,
  fonttitle=\sffamily\bfseries\small,coltitle=ink,title={#1},
  attach title to upper={\par\smallskip}}
\newtcolorbox{achtung}[1][]{enhanced,breakable,colback=white,colframe=white,
  borderline west={3pt}{0pt}{akzent},arc=0pt,left=10pt,right=8pt,top=6pt,bottom=6pt,
  fonttitle=\sffamily\bfseries\small,coltitle=akzent,title={#1},
  attach title to upper={\par\smallskip}}

% --- Code / Quellenverweise ---------------------------------------------
\lstset{basicstyle=\ttfamily\footnotesize,columns=fullflexible,keepspaces=true,
  breaklines=true,frame=leftline,framerule=1.5pt,rulecolor=\color{rule},
  xleftmargin=10pt,framexleftmargin=6pt,showstringspaces=false,
  literate={ä}{{\"a}}1 {ö}{{\"o}}1 {ü}{{\"u}}1 {Ä}{{\"A}}1 {Ö}{{\"O}}1 {Ü}{{\"U}}1 {ß}{{\ss}}1
           {→}{{$\rightarrow$}}1 {–}{{--}}1 {−}{{$-$}}1}
\DeclareUrlCommand\datei{\def\UrlFont{\ttfamily\small}}
\newcommand{\code}[1]{\texttt{#1}}
\newcommand{\quelle}[1]{\hfill{\footnotesize\color{ink3}\textsf{Quelle:} \texttt{#1}}}

% --- Mathematik ----------------------------------------------------------
\newcommand{\Pois}{\operatorname{Pois}}
\newcommand{\clamp}{\operatorname{clamp}}
\newcommand{\sgn}{\operatorname{sgn}}
\newcommand{\E}{\mathbb{E}}

\hypersetup{pdftitle={League Simulator -- Handbuch},pdfauthor={Christoph Schwerdtfeger}}

\begin{document}

\begin{titlepage}
  \raggedright
  \vspace*{3cm}
  {\sffamily\small\color{ink3} 30 Punkte · fussball.csdatascience.de\par}
  \vspace{6pt}
  {\color{akzent}\rule{40pt}{2pt}\par}
  \vspace{14pt}
  {\fontsize{38}{42}\selectfont League Simulator\par}
  \vspace{10pt}
  {\Large\color{ink2} Methodik, Website und Technik\par}
  \vspace{2cm}
  \begin{minipage}{0.8\textwidth}
    \color{ink2}
    Ein Handbuch zu dem System, das für zehn deutsche Fußballligen nach jedem
    Spiel neu ausrechnet, wie die Saison ausgehen könnte: wie das Modell
    rechnet, was die Website zeigt und worauf das alles läuft.
  \end{minipage}
  \vfill
  {\sffamily\small\color{ink3}
  Stand: Repository \texttt{main} @ \texttt{a41890a}, 11.~Oktober 2026 · Saison 2026/27\par}
\end{titlepage}

\tableofcontents

\chapter{Überblick}

Der League Simulator schätzt für zehn deutsche Fußballligen, mit welcher
Wahrscheinlichkeit jede Mannschaft am Saisonende auf welchem Tabellenplatz
steht. Er rechnet nach jedem beendeten Spiel neu und veröffentlicht das
Ergebnis als statische Website unter
\href{https://fussball.csdatascience.de}{fussball.csdatascience.de}, die unter
dem Namen des Blogs \enquote{30~Punkte} läuft.

Abgedeckt sind die Bundesliga, die 2.~Bundesliga, die 3.~Liga, die fünf
Regionalligen (Nord, Nordost, West, SüdWest, Bayern) sowie die
Frauen-Bundesliga und die 2.~Frauen-Bundesliga.

\section*{Das Modell in drei Sätzen}

\begin{enumerate}
  \item Jede Mannschaft hat eine \emph{Stärkeschätzung}, den ELO-Wert. Er steigt,
        wenn sie besser abschneidet, als das Modell erwartet hat, und fällt,
        wenn sie schlechter abschneidet (Kapitel~\ref{kap:elo}).
  \item Aus dem ELO-Unterschied zweier Mannschaften folgt für jede eine
        Torerwartung. Die Tore eines Spiels werden als unabhängige
        Poisson-Zufallszahlen gezogen (Kapitel~\ref{kap:tormodell}).
  \item Der Rest der Saison wird zehntausendmal durchgespielt. Der Anteil der
        Durchläufe, in denen eine Mannschaft auf einem Platz landet, ist die
        angezeigte Wahrscheinlichkeit (Kapitel~\ref{kap:simulation}).
\end{enumerate}

Alles Weitere ist Feinarbeit: Regeln für Auf- und Abstieg zwischen den Ligen
(Kapitel~\ref{kap:regionalliga}), die Herkunft der Startwerte
(Kapitel~\ref{kap:startwerte}) und eine ehrliche Liste dessen, was das Modell
nicht kann (Kapitel~\ref{kap:grenzen}).

\section*{Aufbau dieses Handbuchs}

\begin{description}[font=\sffamily\bfseries\small]
  \item[Teil I, Methodik] Alle Formeln, so wie der Code sie rechnet,
        mit Konstanten, Rechenbeispielen und Quellverweisen.
  \item[Teil II, Website] Die 22 Seiten und jedes Element darauf, von der
        Heatmap bis zum ELO-Verlaufsdiagramm, dazu das Gestaltungssystem.
  \item[Teil III, Technik] Datenfluss, Scheduler, Rust-Simulationsserver,
        Rendering, Container, CI und Betrieb.
\end{description}

\begin{merke}[Zur Genauigkeit]
Die Formeln in diesem Handbuch sind aus dem Quelltext abgelesen und nicht
aus Lehrbüchern übernommen. Wo der Code von der üblichen Darstellung
abweicht, etwa beim Tordifferenz-Multiplikator oder bei der Begrenzung der
ELO-Differenz, gilt der Code. Verweise wie
\datei{league-simulator-rust/src/elo/mod.rs:5} nennen Datei und Zeile auf dem
Stand \texttt{a41890a}. Alle Rechenbeispiele sind nachgerechnet.
\end{merke}

\section*{Die Modellkonstanten}

Zwei Parametergruppen bestimmen jede Prognose. Sie stehen
\emph{ausschließlich} im Rust-Server
(\datei{league-simulator-rust/src/models/mod.rs:124--130}). R sendet einen Wert
nur, wenn eine Liga bewusst abweicht, und das betrifft nur das Tormodell der
Frauen-Ligen (ADR~0002, ADR~0004). Ein Wächtertest im R-Testsatz sorgt dafür,
dass keine dieser Zahlen im R- oder JavaScript-Code oder im Seitentext
auftaucht.

\begin{center}
\begin{tabular}{@{}llll@{}}
\toprule
Symbol & Name im Code & Wert & Gilt für \\
\midrule
$k$   & \code{mod\_factor}      & 20 & alle Ligen \\
$H$   & \code{home\_advantage}  & 40 & alle Ligen \\
$N$   & \code{iterations}       & \num{10000} & alle Ligen \\
$s$   & \code{tore\_slope}      & 0{,}0017854953143549 & Herren-Ligen \\
$c$   & \code{tore\_intercept}  & 1{,}3218390804597700 & Herren-Ligen \\
$s_F$ & \code{tore\_slope}      & 0{,}0024058833 & Frauen-Ligen (82, 1034) \\
$c_F$ & \code{tore\_intercept}  & 1{,}6527603153 & Frauen-Ligen (82, 1034) \\
\bottomrule
\end{tabular}
\end{center}


\part{Methodik}
\chapter{Die Stärkeschätzung: ELO}\label{kap:elo}

Das ELO-System stammt aus dem Schach. Jede Mannschaft trägt eine Zahl, die
ihre Spielstärke schätzt. Nach jedem Spiel wird verglichen, was das Modell
erwartet hat und was passiert ist. Die Differenz, verstärkt durch die Höhe
des Ergebnisses, verschiebt die Zahlen beider Mannschaften um denselben
Betrag in entgegengesetzte Richtungen.

In der Bundesliga liegen die Werte grob zwischen 1\,500 und 2\,100, in den
Regionalligen etwa zwischen 700 und 1\,200. Die Skala hat keinen absoluten
Nullpunkt. Bedeutung haben nur Unterschiede, und das auch nur innerhalb einer
Wechselgemeinschaft, also einer Gruppe von Ligen, zwischen denen Mannschaften
auf- und absteigen (Abschnitt~\ref{sec:wechselgemeinschaft}).

Die Rechnung steht in einer einzigen Funktion,
\code{calculate\_elo\_change}, in
\datei{league-simulator-rust/src/elo/mod.rs:5--32}. Sie übernimmt die Logik
der früheren C++-Implementierung (\datei{SpielNichtSimulieren.cpp}).

\section{Erwartungswert}

Es seien $E_h$ und $E_a$ die ELO-Werte der Heim- und der Gastmannschaft
vor dem Spiel und $H = 40$ der Heimvorteil. Der Code bildet zuerst die
\emph{invertierte} Differenz und begrenzt sie auf $\pm 400$:
\begin{equation}
  \Delta_{\text{inv}} = E_a - E_h - H, \qquad
  \tilde\Delta = \clamp_{[-400,\,400]}\bigl(\Delta_{\text{inv}}\bigr)
              = \min\bigl(\max(\Delta_{\text{inv}}, -400),\, 400\bigr).
\end{equation}
Der Erwartungswert der Heimmannschaft ist die klassische logistische
ELO-Formel mit Basis 10 und Skala 400:
\begin{equation}\label{eq:elo-erwartung}
  p_h = \frac{1}{1 + 10^{\tilde\Delta/400}}.
\end{equation}

\begin{itemize}
  \item Der Heimvorteil wird auf die ELO der Heimmannschaft aufgeschlagen.
        Bei gleich starken Mannschaften erwartet das Modell für die
        Heimmannschaft $p_h = 1/(1+10^{-0{,}1}) \approx 0{,}557$.
  \item Wegen der Begrenzung gilt immer
        $p_h \in [\tfrac{1}{11},\, \tfrac{10}{11}] \approx [0{,}091;\ 0{,}909]$.
        Auch ein krasser Favorit kann bei einem Sieg also noch Punkte gewinnen.
  \item $p_h$ ist ein \emph{erwarteter Score}, bei dem ein Remis als halber
        Sieg zählt. Es ist \emph{nicht} die Siegwahrscheinlichkeit des
        Tormodells (Abschnitt~\ref{sec:skalen}).
\end{itemize}

\section{Tatsächlicher Score und Tordifferenz-Multiplikator}

Das Ergebnis wird auf einen Score abgebildet. Ein Sieg zählt 1, ein Remis
0{,}5 und eine Niederlage 0
(\datei{elo/mod.rs:16--17}):
\begin{equation}
  S = \frac{\sgn(g_h - g_a) + 1}{2} \in \{0,\ \tfrac12,\ 1\}.
\end{equation}
Hohe Siege zählen mehr, aber gedämpft. Der Multiplikator ist die Wurzel der
Tordifferenz, mindestens 1 (\datei{elo/mod.rs:19--20}):
\begin{equation}
  m = \sqrt{\max\bigl(\lvert g_h - g_a\rvert,\ 1\bigr)}.
\end{equation}
Remis und Siege mit einem Tor Unterschied haben $m = 1$. Ein 4:0 hat
$m = 2$ und verschiebt damit doppelt so viel wie ein 1:0.

\section{Das Update}

Mit dem Gewichtungsfaktor $k = 20$ (\code{mod\_factor}) gilt
(\datei{elo/mod.rs:23--27}):
\begin{equation}\label{eq:elo-update}
  \boxed{\;\delta = k \cdot m \cdot (S - p_h), \qquad
  E_h' = E_h + \delta, \qquad E_a' = E_a - \delta\;}
\end{equation}

\begin{description}[font=\sffamily\bfseries\small]
  \item[Nullsumme] Was die eine Mannschaft gewinnt, verliert die andere. In
        einer geschlossenen Liga bleiben Summe und Mittelwert der ELO über die
        ganze Saison konstant. Der Rust-Test
        \code{elo\_is\_conserved\_by\_the\_walk} prüft das, und die Eichung
        neuer Ligen (Abschnitt~\ref{sec:eichung}) nutzt es.
  \item[Keine Rundung] ELO ist durchgehend eine Gleitkommazahl. Die
        Stammdatei speichert volle Genauigkeit, zum Beispiel
        \code{2057.25650898366} für den FC Bayern. Gerundet wird nur bei der
        Anzeige.
  \item[Reihenfolge] Die Updates werden in chronologischer Reihenfolge der
        Anstoßzeiten ausgeführt, ausgehend vom ELO-Wert zum Saisonstart
        (\code{InitialELO} in \datei{RCode/TeamList_2026.csv}). Diese
        Kette heißt im Code \emph{ELO-Walk}.
  \item[Gewertete Spiele] Spiele, die am grünen Tisch gewertet wurden
        (Status AWD oder WO), zählen für die Tabelle, bewegen die ELO aber
        nicht (Issue~\#157). Ein 3:0 am grünen Tisch sagt nichts über die
        Spielstärke aus.
\end{description}

\section{Rechenbeispiel}\label{sec:elo-beispiel}

Das Eröffnungsspiel der Saison 2026/27: FC Bayern ($E_h = 2057{,}26$) gegen
VfB Stuttgart ($E_a = 1805{,}06$).
\begin{align*}
  \Delta_{\text{inv}} &= 1805{,}06 - 2057{,}26 - 40 = -292{,}19
    \quad\text{(innerhalb $\pm 400$)},\\
  p_h &= 1/(1 + 10^{-292{,}19/400}) = 0{,}8432.
\end{align*}

\begin{center}
\begin{tabular}{@{}lccr@{}}
\toprule
Ergebnis & $S$ & $m$ & ELO-Änderung Bayern $\delta$ \\
\midrule
2:1 & 1   & 1          & $20 \cdot 1 \cdot (1 - 0{,}8432) = +3{,}14$ \\
4:0 & 1   & 2          & $+6{,}27$ \\
1:1 & 0{,}5 & 1        & $-6{,}86$ \\
0:2 & 0   & $\sqrt2$   & $-23{,}85$ \\
\bottomrule
\end{tabular}
\end{center}

Stuttgart erhält jeweils den Gegenwert mit umgekehrtem Vorzeichen. Ein Remis
kostet den Favoriten mehr, als ein knapper Sieg ihm einbringt: Das Modell
hatte einen Sieg weitgehend vorausgesetzt.

Eine zweite Gegenprobe steht in ADR~0002. Ein 3:0 bei 1\,500 gegen 1\,480
ergibt $\delta = 20\sqrt3\,(1 - p_h) = +14{,}359$.

\begin{merke}[Wie die ELO-Werte auf die Seite kommen]
Gespeichert wird nur die ELO zum Saisonstart. Den aktuellen Stand und jeden
Zwischenstand rechnet der Rust-Server bei jedem Rendern neu aus allen
gespielten Spielen nach. Das geschieht deterministisch über den Endpunkt
\code{POST /league-details} (ADR~0002). Ligatabelle, Rückblick, Ausblick und
ELO-Verlauf beruhen so alle auf demselben Walk und können sich nicht
widersprechen.
\end{merke}

\chapter{Vom Stärkeunterschied zum Ergebnis: das Tormodell}\label{kap:tormodell}

Für die Prognose reicht ein erwarteter Score nicht. Die Simulation braucht
konkrete Ergebnisse, denn Tabellenplätze hängen auch an Toren und
Tordifferenz. Das Tormodell übersetzt den ELO-Unterschied in zwei
Torerwartungen und zieht daraus Tore.

\section{Torerwartung}

Mit der Differenz einschließlich Heimvorteil
\begin{equation}
  \Delta = E_h + H - E_a
\end{equation}
sind die erwarteten Tore der beiden Mannschaften
(\datei{league-simulator-rust/src/simulation/match_sim.rs:16--21}):
\begin{equation}\label{eq:lambda}
  \boxed{\;\lambda_h = \max\bigl(s\,\Delta + c,\ 0{,}001\bigr), \qquad
          \lambda_a = \max\bigl(-s\,\Delta + c,\ 0{,}001\bigr)\;}
\end{equation}
Für die Herren-Ligen gilt $s = 0{,}0017854953143549$ und
$c = 1{,}32183908045977$, für die Frauen-Ligen $s_F = 0{,}0024058833$ und
$c_F = 1{,}6527603153$.

\begin{itemize}
  \item Das Modell ist \emph{affin} in $\Delta$. Die Differenz wird hier
        nicht begrenzt, anders als beim ELO-Update. Es gibt nur die
        Untergrenze $\lambda \ge 0{,}001$. Sie greift erst ab
        $\lvert\Delta\rvert > (c - 0{,}001)/s$, also ab 739{,}8 ELO-Punkten
        bei den Herren und ab 686{,}6 bei den Frauen.
  \item Solange die Untergrenze nicht greift, ist $\lambda_h + \lambda_a = 2c$.
        Die erwartete Torzahl je Spiel hängt damit nicht von den Mannschaften
        ab: $2c = 2{,}64$ bei den Herren und $3{,}31$ bei den Frauen.
        Kapitel~\ref{kap:grenzen} behandelt die Folgen.
  \item Dieselbe Formel benutzt auch der deterministische Endpunkt
        \code{/league-details}
        (\datei{league_details/mod.rs:82--85}). Simulation und Spieldetails
        rechnen also mit demselben Modell.
\end{itemize}

\section{Ziehen der Tore}

In der Simulation werden für jedes Spiel zwei unabhängige gleichverteilte
Zufallszahlen $u_h, u_a \in [0, 1)$ gezogen. Über die Quantilfunktion der
Poisson-Verteilung werden sie zu Toren
(\datei{match_sim.rs:41--132}):
\begin{equation}
  g = F^{-1}_{\Pois(\lambda)}(u) = \min\bigl\{\,j \in \mathbb{N}_0 :
      P(X \le j) \ge u \,\bigr\}, \qquad X \sim \Pois(\lambda).
\end{equation}
Das entspricht \code{qpois} in R. Für $\lambda < 10$, also in jedem
praktischen Fall, summiert der Code die Wahrscheinlichkeiten direkt auf, mit
der Rekursion $P(X = j) = P(X = j-1)\cdot\lambda/j$. Für größere $\lambda$
sucht er binär über die Verteilungsfunktion aus der Bibliothek \code{statrs}.

Die Tore der beiden Mannschaften sind \emph{unabhängig}. Es gibt keinen
Korrelationsterm wie im Dixon-Coles-Modell. Diese Vereinfachung erklärt die
bekannten Abweichungen bei der Remisquote (Abschnitt~\ref{sec:remis}).

\section{Sieg, Remis, Niederlage in geschlossener Form}

Die Spieldetails auf der Website (Rückblick, Ausblick, Tooltip im
ELO-Verlauf) brauchen keine Zufallszahlen. Der Rust-Server rechnet die
Wahrscheinlichkeiten exakt aus
(\datei{league_details/mod.rs:61--115}). Mit
$\pi_\lambda(j) = e^{-\lambda}\lambda^j/j!$ für $j = 0, \dots, 30$ gilt:
\begin{align}
  P(\text{Heimsieg})     &= \sum_{h=0}^{30}\sum_{a=0}^{h-1} \pi_{\lambda_h}(h)\,\pi_{\lambda_a}(a), \\
  P(\text{Remis})        &= \sum_{j=0}^{30} \pi_{\lambda_h}(j)\,\pi_{\lambda_a}(j), \\
  P(\text{Auswärtssieg}) &= \sum_{a=0}^{30}\sum_{h=0}^{a-1} \pi_{\lambda_h}(h)\,\pi_{\lambda_a}(a).
\end{align}
Die Summen brechen bei 30 Toren je Mannschaft ab und werden nicht
renormiert. Bei den vorkommenden $\lambda$ ist die abgeschnittene Masse
kleiner als $10^{-12}$.

\subsection*{Die Ergebnis-Matrix}

Für den Ausblick berechnet derselbe Code die Wahrscheinlichkeit jedes
Endstands. Mit der Standardeinstellung \code{max\_goals = 6} entsteht ein
$7 \times 7$-Gitter, dessen letzte Zeile und Spalte die Restmasse
\enquote{6+} sammeln:
\begin{equation}
  M_{ij} = \sum_{\substack{h,\,a \le 30\\ \min(h,6) = i,\ \min(a,6) = j}}
           \pi_{\lambda_h}(h)\,\pi_{\lambda_a}(a), \qquad i, j \in \{0, \dots, 6\}.
\end{equation}

\section{Rechenbeispiel}\label{sec:tor-beispiel}

Wieder FC Bayern gegen VfB Stuttgart (Herren-Tormodell):
\begin{align*}
  \Delta &= 2057{,}26 + 40 - 1805{,}06 = 292{,}19,\\
  \lambda_h &= \phantom{-}0{,}0017855 \cdot 292{,}19 + 1{,}32184 = 1{,}844,\\
  \lambda_a &= -0{,}0017855 \cdot 292{,}19 + 1{,}32184 = 0{,}800.
\end{align*}
Daraus folgen $P(\text{Sieg Bayern}) = 62{,}1\,\%$,
$P(\text{Remis}) = 22{,}2\,\%$ und
$P(\text{Sieg Stuttgart}) = 15{,}7\,\%$. Die Methodik-Seite rundet das auf
62, 22 und 16~Prozent und übersetzt es so: \enquote{Liefe dieser Abend
sechsmal, gewönne Stuttgart einmal.}

Ausschnitt der Ergebnis-Matrix in Prozent (Zeilen: Tore Bayern, Spalten:
Tore Stuttgart):

\begin{center}
\begin{tabular}{@{}c|cccc@{}}
     & 0 & 1 & 2 & 3 \\ \midrule
0    & 7{,}1  & 5{,}7  & 2{,}3 & 0{,}6 \\
1    & 13{,}1 & 10{,}5 & 4{,}2 & 1{,}1 \\
2    & 12{,}1 & 9{,}7  & 3{,}9 & 1{,}0 \\
3    & 7{,}4  & 5{,}9  & 2{,}4 & 0{,}6 \\
\end{tabular}
\end{center}
Das wahrscheinlichste Einzelergebnis ist 1:0 mit 13{,}1~Prozent.

\section{Zwei Erwartungen, zwei Skalen}\label{sec:skalen}

\begin{achtung}[Eine Eigenheit, die man kennen sollte]
Das ELO-Update rechnet mit der logistischen Erwartung aus
Gleichung~\eqref{eq:elo-erwartung}, die Prognose mit dem Poisson-Tormodell.
Beide Zahlen messen etwas Ähnliches, sind aber verschieden skaliert. Für
Bayern gegen Stuttgart gilt:
\[
  p_h^{\text{ELO}} = 0{,}843
  \qquad\text{gegenüber}\qquad
  P(\text{Sieg}) + \tfrac12 P(\text{Remis}) = 0{,}621 + 0{,}111 = 0{,}732.
\]
Das ELO-Update erwartet vom Favoriten also mehr als das Tormodell. Ein
knapper Bayern-Sieg bringt deshalb nur $+3{,}1$ Punkte, obwohl das
Tormodell ihn für keineswegs sicher hält. Die Abweichung ist
systematisch und stammt aus der Übernahme des ursprünglichen C++-Codes. Ob
eine einheitliche Skala die Prognosen verbessert, ist nicht untersucht. ADR~0002 beschreibt denselben Skaleneffekt bei der
Kalibrierung des Heimvorteils.
\end{achtung}

\chapter{Zehntausend Saisons: Simulation und Monte Carlo}\label{kap:simulation}

\section{Eine simulierte Saison}

Eine Iteration spielt die ganze Saison durch
(\datei{league-simulator-rust/src/simulation/season.rs:10--62}). Sie beginnt
jedes Mal bei den ELO-Werten zum Saisonstart und geht den Spielplan in
Listenreihenfolge durch. Für jedes Spiel gibt es drei Fälle:

\begin{enumerate}
  \item \textbf{Das Ergebnis fehlt.} Die Tore werden mit dem Tormodell
        gezogen und eingetragen. \emph{Danach wird die ELO aktualisiert}, als
        wäre das Spiel wirklich so ausgegangen.
  \item \textbf{Das Ergebnis liegt vor.} Es zählt für die Tabelle, und die ELO
        wird mit dem echten Ergebnis aktualisiert.
  \item \textbf{Das Ergebnis liegt vor, das Spiel ist gewertet} (Kennzeichen
        \code{elo\_neutral}). Es zählt für die Tabelle, die ELO bleibt
        unverändert.
\end{enumerate}

Die ELO ist innerhalb einer simulierten Saison also \emph{dynamisch}. Eine
Mannschaft, die in einem Durchlauf eine Siegesserie hinlegt, geht stärker in
die folgenden simulierten Spiele. Das bildet ab, dass sich auch die echte
Stärkeschätzung im Lauf der Saison bewegen wird, und verbreitert die
Verteilung der Endplätze gegenüber einem Modell mit festen Stärken.

\subsection*{Reihenfolge der Spiele}

Die R-Seite sortiert den Spielplan chronologisch, bevor er an Rust geht
(\datei{RCode/transform_data.R:271--291}): nach Anstoßzeit, bei Gleichstand
in der Reihenfolge der API, Spiele ohne Termin ans Ende. Gespielte Spiele
werden so in jeder Iteration in derselben Reihenfolge nachgespielt und führen
zu denselben ELO-Werten. Die simulierten Spiele folgen danach in der Reihenfolge
des Spielplans.

\section{Statusgruppen}

Die Daten kommen von API-Football. Jedes Spiel trägt einen Statuscode, und
das Modell fasst die Codes in Gruppen zusammen
(\datei{RCode/league_details.R:56--68},
\datei{RCode/transform_data.R:357--374}):

\begin{center}
\small
\begin{tabularx}{\textwidth}{@{}llXcc@{}}
\toprule
Gruppe & Codes & In der Simulation & ELO-Walk & Tabelle \\
\midrule
beendet    & FT, AET, PEN & Ergebnis fest & ja & ja \\
gewertet   & AWD, WO      & Ergebnis fest, \code{elo\_neutral} & nein & ja \\
live       & 1H, HT, 2H, ET, BT, P, SUSP, INT, LIVE & wird simuliert & nein & nein \\
verschoben & PST, CANC, TBD, ABD & wird simuliert & nein & nein \\
offen      & alle übrigen, vor allem NS & wird simuliert & nein & nein \\
\bottomrule
\end{tabularx}
\end{center}

Laufende Spiele fließen erst mit dem Abpfiff ein. Der Zwischenstand bleibt
unberücksichtigt, und die Website sagt das in der Live-Sektion ausdrücklich.
Simuliert werden nur Hauptrundenspiele. Runden mit Bezeichnungen wie
\enquote{Relegation}, \enquote{Play-off} oder \enquote{Final} filtert
\datei{RCode/round_filter.R} heraus.

\section{Die Abschlusstabelle und ihre Reihenfolge}\label{sec:tiebreak}

Am Ende jeder Iteration entsteht eine Tabelle mit drei Punkten für einen
Sieg, einem für ein Remis und keinem für eine Niederlage. Sortiert wird
stabil (\datei{season.rs:92--168}):

\begin{enumerate}
  \item Punkte, absteigend,
  \item Tordifferenz, absteigend,
  \item erzielte Tore, absteigend,
  \item bei völligem Gleichstand die Eingabereihenfolge der Mannschaften.
\end{enumerate}

Die Eingabereihenfolge ist alphabetisch nach Kürzel
(\datei{transform_data.R:383--387}). Bei völligem Gleichstand gewinnt also
systematisch das alphabetisch frühere Kürzel. Einen direkten Vergleich wertet
das Modell nicht aus. Bei zehntausend Durchläufen und ganzzahligen Toren ist
ein Gleichstand in allen drei Kriterien selten, kommt aber vor.

Die Engine kann Startwerte für Punkte und Tore übernehmen
(\code{adj\_points}, \code{adj\_goals} und weitere). Im Hauptlauf sind sie
null. Reale Punktabzüge, etwa nach Lizenzverstößen, werden nicht modelliert,
weil es dafür keine Datenquelle gibt. Genutzt werden die Startwerte nur für
den Malus-Lauf der Zweitvertretungen (Abschnitt~\ref{sec:malus}).

\section{Monte Carlo}

Der Rust-Server wiederholt die Saison $N = \num{10000}$-mal
(\datei{league-simulator-rust/src/monte_carlo/mod.rs}). Es sei $C_{t,p}$ die
Zahl der Iterationen, in denen Mannschaft $t$ auf Platz $p$ landet. Die
angezeigte Wahrscheinlichkeit ist der Anteil:
\begin{equation}\label{eq:mc}
  \boxed{\;\hat P(t \text{ auf Platz } p) = \frac{C_{t,p}}{N}\;}
\end{equation}
Jede Zeile und jede Spalte dieser Matrix summiert sich zu 1. Die Zeilen
sortiert der Server nach dem erwarteten Platz
$\bar p_t = \sum_p p\,\hat P_{t,p}$, aufsteigend. Diese Matrix ist die
Heatmap auf jeder Ligaseite.

\begin{description}[font=\sffamily\bfseries\small]
  \item[Zufallszahlen] Jede Iteration erhält einen eigenen 64-Bit-Startwert
        aus dem Thread-Zufallsgenerator des Betriebssystems und erzeugt daraus
        einen eigenen Generator (\code{StdRng::seed\_from\_u64}). Zwei Läufe
        mit denselben Daten liefern deshalb leicht unterschiedliche Matrizen.
        Eine deterministische Variante mit festem Startwert gibt es nur für
        Tests.
  \item[Parallelität] Die Iterationen laufen über \code{rayon} parallel
        (\code{par\_iter().fold().reduce()}). Jeder Thread zählt in eigene
        Zähler, die am Ende addiert werden. Ganzzahlige Addition ist
        kommutativ, die Reihenfolge der Threads ändert das Ergebnis also nicht.
  \item[Laufzeit] Rund 42~ms für \num{10000} Iterationen einer Liga mit 18
        Mannschaften und vollem Restspielplan, gemessen am 7.~September 2026.
        Alle zehn Ligen zusammen brauchen etwa eine halbe Sekunde.
\end{description}

\subsection*{Wie genau sind die Zahlen?}

Der Standardfehler eines geschätzten Anteils ist
\begin{equation}
  \operatorname{SE}(\hat P) = \sqrt{\frac{\hat P (1 - \hat P)}{N}} \le
  \sqrt{\frac{0{,}25}{\num{10000}}} = 0{,}5\ \text{Prozentpunkte}.
\end{equation}
Die Methodik-Seite fasst das allgemeinverständlich zusammen: 35 gegen 37
Prozent heißt gleichauf. Ein Häkchen (\checkmark) bedeutet, dass in keiner
der zehntausend Simulationen etwas anderes herauskam. Das ist etwas anderes
als \enquote{rechnerisch sicher}.

\section{Von der Matrix zu den Spalten der Website}

Die Spalten neben der Heatmap, etwa \enquote{Meister},
\enquote{Champions League} oder \enquote{Abstieg}, summieren Plätze
(\code{groupResultsDF}, \datei{RCode/render_helpers.R:49--77}):
\begin{equation}
  P(t \in \text{Gruppe } [l, u]) = \sum_{p=l}^{u} \hat P_{t,p}.
\end{equation}
Negative Grenzen zählen von unten: $-1$ ist der letzte Platz. Sie werden erst
zur Renderzeit gegen die Ligagröße aufgelöst, weil Ligen ihre Größe ändern
können (die Frauen-Bundesliga wuchs 2025 von 12 auf 14 Mannschaften).
Tabelle~\ref{tab:panels} in Kapitel~\ref{kap:ligaseite} listet alle Gruppen
je Liga.

\subsection*{Zusatzzählung für die 3.~Liga}

Für die Regionalligen ist nicht nur wichtig, wie wahrscheinlich ein
Drittligist absteigt, sondern auch \emph{wohin}. Er fällt in die Regionalliga
seiner Region zurück. Deshalb zählt der Server für die 3.~Liga zusätzlich,
wie viele der vier Absteiger aus jeder Staffelregion kommen. Es sei $G_{g,j}$
die Zahl der Iterationen mit genau $j$ Absteigern aus Region $g$
(\datei{monte_carlo/mod.rs:133--147}). Auf diesen Zählungen baut die
Abstiegskopplung auf (Abschnitt~\ref{sec:kopplung}).

\chapter{Regionalligen: Kopplung, Aufstieg, Malus}\label{kap:regionalliga}

In den Regionalligen hängt an einem Tabellenplatz oft nicht allein, ob man
auf- oder absteigt. Wie viele Mannschaften eine Staffel abgeben muss, hängt
davon ab, wie viele Drittligisten in sie zurückfallen. Wer aufsteigt, hängt
für zwei Staffeln an Aufstiegsspielen gegeneinander. Die Rechtsgrundlagen
für 2026/27 stehen in \datei{docs/abstieg_aufstieg_RL_2026_2027.md}, die
Modellannahmen in \datei{docs/modellannahmen-rl-kopplung.md}. Die
Entscheidungen dazu halten ADR~0006 und ADR~0008 fest.

\section{Die Abstiegskopplung an die 3.~Liga}\label{sec:kopplung}

\subsection*{Verteilung der Drittliga-Absteiger je Staffel}

Jeder Drittligist ist über die Spalte \code{Region} der Stammdatei einer
Staffel zugeordnet. Aus der Zusatzzählung des Monte-Carlo-Laufs
(Kapitel~\ref{kap:simulation}) folgt direkt die Verteilung der Zahl der
Absteiger je Staffel. Sie ist ausgezählt und nicht genähert
(\datei{RCode/rl_abstiegskopplung.R:116--181}):
\begin{equation}
  q_g(j) = P(\text{genau } j \text{ Drittliga-Absteiger in Staffel } g) = \frac{G_{g,j}}{N}.
\end{equation}
Ausgezählt und nicht etwa aus den einzelnen Abstiegswahrscheinlichkeiten
zusammengesetzt wird, weil die Abstiegsplätze negativ korreliert sind:
Steigt ein Verein aus dem Norden ab, ist ein Platz weniger für alle anderen
frei. Der Code prüft zwei Invarianten. Jede Zeile summiert sich zu $N$, und
die erwartete Gesamtzahl $\sum_g \sum_j j\,q_g(j)$ ist exakt $K = 4$.

\subsection*{Abstiegsplätze je Staffel}

Wie viele Plätze bei $j$ Drittliga-Absteigern zum Abstieg führen, regeln die
Spielordnungen der Regionalverbände (\datei{rl_abstiegskopplung.R:65--86}):

\begin{center}
\begin{tabular}{@{}llp{6.2cm}@{}}
\toprule
Staffel & Abstiegsplätze $n_g(j)$ & Grundlage \\
\midrule
Nord     & $3 + j$             & NFV-SpO §\,6 Abs.~3/4 \\
Nordost  & $\min(1 + j,\ 2)$   & NOFV A.\ Nr.~5, Deckel 2 als Modellannahme \\
West     & $4$                 & WDFV, entkoppelt \\
SüdWest  & $\min(3 + j,\ 5)$   & RLSW-SpO §\,47 Nr.~1/2 \\
Bayern   & $2$                 & BFV, dazu Relegation (s.\,u.) \\
\bottomrule
\end{tabular}
\end{center}

\subsection*{Platzgewichte und Abstiegswahrscheinlichkeit}

Für eine Staffel mit $n$ Mannschaften ist der $d$-letzte Platz genau dann ein
Abstiegsplatz, wenn $n_g(j) \ge d$. Sein Gewicht ist die Wahrscheinlichkeit
dafür (\datei{rl_abstiegskopplung.R:339--345}):
\begin{equation}
  w_g(n - d + 1) = \sum_{j \,:\, n_g(j) \ge d} q_g(j), \qquad d = 1, \dots, n.
\end{equation}
Die Abstiegswahrscheinlichkeit einer Mannschaft verbindet dieses Gewicht mit
ihrer eigenen Platzverteilung aus dem Hauptlauf der Staffel
(\datei{rl_abstiegskopplung.R:371--392}):
\begin{equation}\label{eq:rl-abstieg}
  \boxed{\;P(t \text{ steigt ab}) = \sum_{p=1}^{n} \hat P_{t,p}\; w_g(p)\;}
\end{equation}
Das Produkt ist exakt und keine Näherung. 3.~Liga und Regionalliga sind
getrennte Wettbewerbe, ihre Ergebnisse also unabhängig.

\paragraph{Nord: die zweite Kopplung.} Steigt der Nord-Meister auf, muss die
Staffel einen Absteiger weniger stellen. Mit $p_M$ als Summe der
Aufstiegswahrscheinlichkeiten der Staffel mischt der Code beide Fälle
(\datei{RCode/update_all_leagues_loop.R:775--780}):
\begin{equation}
  w^{\text{Nord}} = \clamp_{[0,1]}\Bigl((1 - p_M)\, w[n_g(j)] + p_M\, w[\max(n_g(j) - 1,\ 0)]\Bigr).
\end{equation}
Hier handelt es sich ausdrücklich um eine Näherung: Meisterschaft und
Abstieg stammen aus derselben Simulation der Staffel und sind nicht
unabhängig.

\paragraph{Bayern: Relegation.} Die beiden Letzten steigen direkt ab. Die
Plätze $n-3$ und $n-2$ spielen eine Relegation gegen die Bayernligen, die das
Modell nicht auflöst. Die Website weist diese Plätze als eigene Spalte
\enquote{Relegation} aus:
$P(\text{Relegation}) = \hat P_{t,n-3} + \hat P_{t,n-2}$.

\section{Wer aufsteigt}

Vier Regionalligisten steigen in die 3.~Liga auf. West und SüdWest stellen
ihren Aufsteiger dauerhaft direkt. Der dritte direkte Platz rotiert unter
Nord, Nordost und Bayern. Die beiden übrigen dieser drei ermitteln den
vierten Aufsteiger in Hin- und Rückspiel. Wer in welcher Saison an der Reihe
ist, steht als Daten in \code{AUFSTIEGSROTATION}
(\datei{RCode/rl_aufstieg.R:53}). Für 2026/27 steigt Nordost direkt auf,
\textbf{Nord und Bayern} spielen gegeneinander. Für eine unbekannte Saison
bricht der Code ab, statt still mit dem Vorjahreswert weiterzurechnen.

\subsection*{Direktaufsteiger}

\begin{equation}
  P(\text{Aufstieg}) = P(\text{Meister}) = \hat P_{t,1}.
\end{equation}

\subsection*{Aufstiegsspiele}

Für die beiden Playoff-Staffeln summiert der Code über alle möglichen
Meisterpaarungen. Weil die Staffeln unabhängig simuliert werden, ist auch
diese Doppelsumme exakt (\datei{rl_aufstieg.R:120--190}):
\begin{equation}\label{eq:rl-aufstieg}
  \boxed{\;P(X \text{ steigt auf}) = P_M(X) \cdot \sum_{Y} P_M(Y)\cdot P(X \text{ setzt sich gegen } Y \text{ durch})\;}
\end{equation}
Dabei ist $P_M$ die Meisterwahrscheinlichkeit aus dem Hauptlauf der
Staffel, demselben, aus dem auch die Spalte \enquote{Meister} der
Aufstiegsseite stammt, und $Y$ läuft über die Mannschaften der anderen
Staffel. Einen Malus-Lauf gibt es für die Regionalligen nicht
(Abschnitt~\ref{sec:malus}). Weil beide Größen aus einem Lauf kommen, gilt
immer $P(\text{Aufstieg}) \le P(\text{Meister})$.

\paragraph{Ein Duell über zwei Spiele.} Die Quote
$P(X \text{ setzt sich gegen } Y \text{ durch})$ rechnet
\datei{RCode/aufstiegsspiele.R:43--117} analytisch aus:

\begin{enumerate}
  \item Die Torerwartungen beider Spiele liefert der Endpunkt
        \code{/match-preview}, mit der \emph{aktuellen} ELO beider
        Mannschaften. Gewertete Spiele bewegen sie auch hier nicht
        (\code{rl\_aktuelle\_elo} sendet \code{elo\_neutral}, Issue~\#293).
        Im Hinspiel hat Nord Heimrecht, im Rückspiel Bayern.
  \item Die Tordifferenz eines Spiels ist verteilt als
        \[
          P(D = d) = \sum_{i} \pi_{\lambda_X}(i)\,\pi_{\lambda_Y}(i - d),
        \]
        abgeschnitten bei 15 Toren je Mannschaft.
  \item Die Gesamtdifferenz ist die Faltung $D_1 * D_2$ der beiden Spiele.
        Eine Auswärtstorregel gibt es nicht mehr.
  \item Bei Gleichstand folgt eine Verlängerung mit den \emph{neutralen}
        Raten $\lambda^{\text{neutral}} = (\lambda^{\text{hin}} + \lambda^{\text{rück}})/2$,
        auf 30 Minuten skaliert: $\lambda_{\text{ev}} = \lambda^{\text{neutral}}/3$.
        Ist es danach immer noch gleich, gilt das Elfmeterschießen als
        50:50.
\end{enumerate}
\begin{equation}
  P(X) = P(D_1 + D_2 > 0) + P(D_1 + D_2 = 0)\,\Bigl[P(D_{\text{ev}} > 0) + \tfrac12\,P(D_{\text{ev}} = 0)\Bigr].
\end{equation}
Im affinen Tormodell hebt sich der Heimvorteil über Hin- und Rückspiel
exakt auf. Wer zuerst Heimrecht hat, spielt für die Quote keine Rolle.

\paragraph{Siegquote.} Die Aufstiegsseite zeigt für die Playoff-Staffeln
zusätzlich die Siegquote
\[
  \frac{P(\text{Aufstieg})}{P(\text{Meister})},
\]
also die Chance, die Aufstiegsspiele zu gewinnen, wenn man Meister wird. Sie
ist die über die möglichen Gegner gemittelte Siegchance aus
Gleichung~\eqref{eq:rl-aufstieg} und deshalb nie größer als 100~Prozent. Ist
$P(\text{Meister}) = 0$, bleibt die Zelle leer.

\section{Aufstiegsberechtigung und Aufstiegssperre}\label{sec:malus}

Zweite Mannschaften von Profivereinen dürfen nicht in jede Liga aufsteigen.
Nach §\,55b Nr.~3.1 der DFB-Spielordnung gilt (ADR~0008, ADR~0009):

\begin{itemize}
  \item In der 3.~Liga und der 2.~Frauen-Bundesliga sind Zweitvertretungen
        grundsätzlich nicht aufstiegsberechtigt.
  \item In der Regionalliga ist eine Zweitvertretung nur dann nicht
        aufstiegsberechtigt, wenn die erste Mannschaft im kommenden Spieljahr
        in der 3.~Liga spielt (\emph{Aufstiegssperre}).
\end{itemize}

Dazu kommt der \emph{Aufstiegsverzicht}: Ein Verein kann für eine Saison auf
den Aufstieg verzichten (TeamList-Spalte \code{Aufstiegsverzicht}).

\paragraph{Masken statt zweitem Lauf.} Das Modell simuliert jede Liga
\emph{einmal}. Die Engine wertet dieselbe Abschlusstabelle zusätzlich unter
einer oder mehreren \emph{Berechtigungsmasken} aus und zählt, auf welchem Rang
\emph{unter den Berechtigten} jedes Team steht. Ein nicht berechtigter Meister
vererbt seinen Aufstieg an den nächsten Berechtigten. Wer nicht berechtigt ist,
steht in der TeamList (\code{Promotion}, \code{Aufstiegsverzicht}) und wird
nicht aus dem Kürzel erraten. Heatmap und Abstieg kommen unverändert aus dem
Tabellenlauf. Bis Oktober 2026 rechnete das Modell stattdessen einen zweiten
Lauf, in dem Gesperrte mit $-50$ Punkten starteten; die Maske ersetzt ihn
(ADR~0009).

\paragraph{Dynamische Sperre durch den Abstieg der Erstmannschaft.}
Spielt die Erstmannschaft einer Regionalliga-Reserve heute in der
2.~Bundesliga, steht erst am Saisonende fest, ob die Reserve gesperrt ist: Die
Erstmannschaft muss direkt absteigen oder die Relegation (Platz~16 gegen den
berechtigten Dritten der 3.~Liga) verlieren. Das Modell rechnet das exakt als
Mischung
\begin{equation}
  P(\text{Aufstieg}) = \sum_S P(S)\, M_S ,
\end{equation}
wobei $S$ die \emph{Menge} der gesperrten Reserven, $P(S)$ ihre gemeinsame
Verteilung und $M_S$ das Staffelergebnis unter der Maske \enquote{$S$ gesperrt}
ist. Das ist exakt, weil 2.~Bundesliga, 3.~Liga und die Staffeln kein Spiel
teilen. Die gemeinsame Verteilung zählt die Engine in der 2.~Bundesliga
(höchstens zwei Teams steigen direkt ab, die Reserven verschiedener Vereine
sind also gekoppelt), die Relegation des Sechzehnten löst R mit der
Gewinnwahrscheinlichkeit des Relegationsduells (zwei Spiele, Verlängerung mit
Heimvorteil des Rückspiel-Gastgebers) auf. Nord und
Bayern hängen an derselben 2.~Bundesliga; die Doppelsumme der Aufstiegsspiele
läuft deshalb über die gemeinsamen Gewichte beider Staffeln. Die Aufstiegsseite
nennt je betroffener Reserve die Wahrscheinlichkeit, dass sie gesperrt ist; die
Aufstiegszone an Platz~1 der Ligatabelle sinkt um die Wahrscheinlichkeit
eines gesperrten Meisters.

Welche Erstmannschaft zu welcher Reserve gehört, steht in den TeamList-Spalten
\code{Erstvertretung\_TeamID} und \code{Erstvertretung\_ShortText}; die API
liefert diese Zuordnung nicht.

\chapter{Woher die Startwerte kommen}\label{kap:startwerte}

Das Modell braucht für jede Mannschaft einen ELO-Wert zum Saisonstart. Er
steht in der Stammdatei \datei{RCode/TeamList_<Saison>.csv}, die als
gepflegtes Stammdatenblatt gilt (ADR~0007). Die Datei hat die Spalten
\code{TeamID}, \code{ShortText}, \code{Promotion}, \code{InitialELO},
\code{League}, \code{Region} und \code{Name}.

\section{Saisonwechsel}

Zum Saisonwechsel rechnet das Werkzeug
\datei{scripts/season_transition.R} die ELO am Ende der alten Saison aus.
Dazu ruft es denselben Endpunkt \code{POST /league-details} auf, den auch die
Website benutzt, je Liga mit deren Tormodell. Einen eigenen ELO-Walk in R
gibt es seit September 2026 nicht mehr (Issue~\#146). Der frühere hatte mit
einem Heimvorteil von 100 statt 40 gerechnet.

\begin{description}[font=\sffamily\bfseries\small]
  \item[Übertrag 1:1] Die Start-ELO der neuen Saison ist die End-ELO der
        alten, ohne Rückführung zur Mitte und unabhängig von Auf- oder
        Abstieg. Ein Aufsteiger nimmt seine Stärke mit.
  \item[Neue Mannschaften] Wer in keiner beobachteten Liga gespielt hat,
        bekommt einen Basiswert: in der 3.~Liga den Mittelwert der End-ELO der
        vier Absteiger der Vorsaison (Ersatzwert 1\,046), in der
        2.~Bundesliga 1\,350, in der Bundesliga 1\,500, sonst 1\,200
        (\datei{RCode/elo_aggregation.R:404--550}).
\end{description}

\section{Eichung für Ligen ohne Historie}\label{sec:eichung}

Mit der Erweiterung auf zehn Ligen im September 2026 kamen sieben Ligen ohne
eigene ELO-Historie dazu: die fünf Regionalligen und die beiden
Frauen-Ligen. Ihre Startwerte hat das einmalige Offline-Skript
\datei{scripts/calibrate_historical_elo.R} geeicht (ADR~0003):

\begin{enumerate}
  \item \textbf{Walk je Liga} über die verfügbare Historie, mit den
        Startwerten 1\,156 (3.~Liga), 920 (Regionalligen) und 1\,500 (Frauen).
        Die End-ELO einer Saison ist der Start der nächsten. Neue Mannschaften
        beginnen beim Startwert. Gewertet werden nur reguläre Spiele der
        Hauptrunde in chronologischer Reihenfolge.
  \item \textbf{Regionalligen an die 3.~Liga.} Der ganze Block wird
        verschoben, und zwar um die mittlere ELO-Differenz zwischen
        Drittliga-Absteigern und Regionalliga-Aufsteigern über die Wechsel
        2020--2025:
        \[
          \Delta_{RL} = \overline{E}(\text{Absteiger 3.~Liga}) - \overline{E}(\text{Aufsteiger RL}).
        \]
  \item \textbf{Staffeln gegeneinander.} Jede Staffel wird um die Hälfte der
        Abweichung verschoben, die ihre Aufsteiger in der 3.~Liga gegenüber
        dem Mittel aller Staffeln zeigen. Die Dämpfung auf die Hälfte
        berücksichtigt, dass je Staffel nur zwei bis sechs Aufsteiger
        vorliegen.
  \item \textbf{2.~Frauen-Bundesliga an die Frauen-Bundesliga}, analog zu
        Schritt~2.
  \item \textbf{Frauen an die Herren}, per Konvention: Die Frauen-Ligen werden
        so verschoben, dass der Mittelwert der Frauen-Bundesliga dem der
        Herren-Bundesliga entspricht (1\,682). Das ist ausdrücklich keine
        Messung, denn zwischen beiden findet kein Austausch statt.
\end{enumerate}

Alle Verschiebungen sind additiv, die Abstände innerhalb einer Liga bleiben
erhalten. Weil der Walk die Summe der ELO erhält, ist keine Iteration nötig.

\section{Ein Tormodell je Wechselgemeinschaft}\label{sec:wechselgemeinschaft}

Die ELO ist das Einzige, was eine Mannschaft über eine Ligagrenze mitnimmt.
Nur wenn Ligen, die Mannschaften austauschen, dasselbe Tormodell benutzen,
bedeutet ein ELO-Wert diesseits und jenseits der Grenze dasselbe. Maßgeblich
ist deshalb nicht die einzelne Liga, sondern die \emph{Wechselgemeinschaft}
(ADR~0004):

\begin{itemize}
  \item \textbf{Herren} (Ligen 78, 79, 80, 83--87) tauschen untereinander
        Mannschaften aus und teilen ein Tormodell.
  \item \textbf{Frauen} (82, 1034) sind untereinander verbunden, nie mit den
        Herren, und haben daher ein eigenes.
\end{itemize}

Die Frauen-Werte $s_F, c_F$ sind per Maximum-Likelihood an 1\,917 Spielen
geschätzt und außerhalb der Stichprobe validiert. Der Brier-Score sinkt um
2{,}96~\%, der Log-Loss um 2{,}66~\%, und das neue Modell ist in allen fünf
Testsaisons besser
(\datei{docs/reports/2026-09-05-frauen-tormodell.md}). Ein eigener
Heimvorteil für die Frauen hat sich nicht als signifikant erwiesen
(Schätzwert 28, $p = 0{,}23$). Es bleibt bei 40.

\subsection*{Wie der Heimvorteil zu 40 kam}

Bis September 2026 rechnete das Modell mit 65 ELO-Punkten. Die
Neukalibrierung integriert über eine realistische Verteilung von
ELO-Unterschieden und trifft mit 40 die gemessenen Anteile der Bundesliga
2024 und 2025:

\begin{center}
\begin{tabular}{@{}lccc@{}}
\toprule
& Heimsieg & Remis & Auswärtssieg \\
\midrule
Modell mit $H = 40$ & 41{,}0\,\% & 24{,}4\,\% & 34{,}7\,\% \\
gemessen            & 41{,}2\,\% & 25{,}0\,\% & 33{,}8\,\% \\
\bottomrule
\end{tabular}
\end{center}

Auf die Saisonprognosen wirkt die Änderung mit weniger als einem
Prozentpunkt (ADR~0002, Nachtrag;
\datei{docs/reports/2026-09-05-heimvorteil-vorher-nachher.md}).

\chapter{Was das Modell kann und was nicht}\label{kap:grenzen}

Wahrscheinlichkeiten sind keine Vorhersagen. Wenn das Modell einer Mannschaft
30~Prozent für den Abstieg gibt und sie bleibt drin, hat es sich nicht
geirrt. Ob ein Modell gut ist, zeigt sich erst über viele Prognosen: Von
allen Ereignissen mit 30~Prozent sollten etwa drei von zehn eintreten. Dieses
Kapitel sammelt die Grenzen, die das Projekt kennt und dokumentiert hat, und
die Vereinfachungen, die sich aus dem Code ergeben.

\section{Dokumentierte Schwächen}

\subsection{Zu wenige Tore bei den Herren}

Die erwartete Torzahl je Spiel ist im Tormodell konstant $2c$. Für die
Herren sind das 2{,}64 Tore, gemessen wurden in der Bundesliga 3{,}18, rund
ein Fünftel mehr. Eine Korrektur müsste alle acht Herren-Ligen gemeinsam
treffen (Abschnitt~\ref{sec:wechselgemeinschaft}). Sie ist für das Projekt
\enquote{Prognosequalität} im Sommer 2027 vorgesehen
(\datei{docs/reports/2026-09-05-liga-empirie-zehn-ligen.md}).

\subsection{Remis: die Form der Verteilung stimmt nicht überall}\label{sec:remis}

Zwei unabhängige Poisson-Verteilungen können eine bestimmte Remisquote nicht
überschreiten. Am höchsten ist sie bei gleich starken Mannschaften ohne
Heimvorteil:
\begin{equation}
  P(\text{Remis}) \le \sum_{j \ge 0} \pi_c(j)^2 =
  \begin{cases} 26{,}1\,\% & \text{Herren } (c = 1{,}322)\\ 23{,}0\,\% & \text{Frauen } (c_F = 1{,}653).\end{cases}
\end{equation}

\begin{itemize}
  \item In der \textbf{2.~Bundesliga} lag die beobachtete Remisquote bei
        26{,}1~\%, praktisch auf dieser Obergrenze. Bei realer Streuung der
        Spielstärken kann das Modell sie nicht erreichen. Es erzeugt dort
        \emph{zu wenige} Remis.
  \item Bei den \textbf{Frauen} liegt das Modell dagegen \emph{zu hoch}:
        19{,}9~\% gegenüber beobachteten 16{,}2~\%, also 3{,}7 Prozentpunkte
        zu viel.
\end{itemize}

Die gemeinsame Ursache ist die Form der Verteilung. Die Gesamttore streuen
stärker, als Poisson erlaubt (Überdispersion 1{,}20 statt 1{,}00). Die
Verteilung der Tordifferenz ist in der Mitte dünner und an den Rändern
dicker als im Modell. Neue Parameterwerte helfen dagegen nicht. Nötig wäre
ein korreliertes Tormodell wie das von Dixon und Coles.


\subsection{Die Streuung der ELO ist kein freier Parameter}

Bei $k = 20$ und rund 34 Spielen je Saison pendelt sich die
Standardabweichung der ELO innerhalb einer Liga bei etwa 95 bis 100 ein. Wo
die echten Stärkeunterschiede größer sind, etwa in der Frauen-Bundesliga,
bildet das Modell sie zu schwach ab. Rechnerisch wären dort 458 nötig, der
Walk erreicht 240 bis 257 (ADR~0003).

\subsection{Annahmen der Regionalliga-Kopplung}

Die Kopplung beruht auf dokumentierten Annahmen
(\datei{docs/modellannahmen-rl-kopplung.md}, §\,5):

\begin{itemize}
  \item Die Zahl der Absteiger aus Nordost ist bei 2 gedeckelt.
  \item Die Nord-Meisterkopplung ist eine Näherung.
  \item West ist entkoppelt.
  \item Die Bayern-Relegation wird nicht aufgelöst.
  \item Die Verlängerung wird mit einem Drittel der Torerwartung gerechnet.
        Empirisch fallen in einer Verlängerung eher 25 bis 30~Prozent der
        Tore eines regulären Spiels.
  \item Sonderfälle wie Lizenzverweigerung, Verzicht oder Zwangsabstieg
        sind ausgelassen.
\end{itemize}

\subsection{Abstieg eines Zweitligisten und Aufstieg seiner Reserve}

Der gleichzeitige Abstieg eines Zweitligisten in die 3.~Liga und Aufstieg
seiner Zweitvertretung aus der Regionalliga wird bisher nicht
berücksichtigt. Nach §\,55b Nr.~3.1 der DFB-Spielordnung entfiele das
Aufstiegsrecht der Zweitvertretung in diesem Fall, weil ihr Verein im
kommenden Spieljahr schon mit einer Mannschaft in der 3.~Liga spielt. Dann
wäre der Nächstplatzierte ihrer Staffel aufstiegsberechtigt (ADR~0008).

Das Modell rechnet die Regionalligen ohne Malus-Lauf
(Abschnitt~\ref{sec:malus}). Um den Fall abzubilden, müsste es die
Simulation der 2.~Bundesliga mit der der Regionalliga koppeln. Die
Größenordnung ist klein: Nach einer Abschätzung vom 17.~Oktober 2026 liegt
die Wahrscheinlichkeit bei rund 1~Prozent, und praktisch betrifft sie nur
Hannover~96~II in der Regionalliga Nord.

\subsection{Keine Prognose-Historie}

Gespeichert wird nichts außer den Startwerten, und Monte Carlo ist nicht
deterministisch. Wie sich etwa die Meisterwahrscheinlichkeit im Lauf der
Saison entwickelt hat, lässt sich deshalb nicht nachträglich rekonstruieren
(ADR~0002). Der ELO-Verlauf dagegen ist deterministisch und wird bei jedem
Rendern neu berechnet.

\section{Vereinfachungen, die sich aus dem Code ergeben}

Die folgenden Punkte stehen in keinem Dokument als Schwäche, ergeben sich
aber beim Lesen des Codes. Einige sind bewusste Vereinfachungen, andere
könnten Kandidaten für Issues sein.

\begin{description}[font=\sffamily\bfseries\small,style=nextline]
  \item[Kein direkter Vergleich]
        Gleichstand nach Punkten, Tordifferenz und Toren entscheidet in der
        Simulation das alphabetisch frühere Kürzel (Abschnitt~\ref{sec:tiebreak}).
  \item[Keine Punktabzüge]
        Reale Abzüge nach Sportgerichtsurteilen oder Lizenzentscheidungen
        gehen nicht ein.
  \item[Abgesagte Spiele werden simuliert]
        Ein Spiel mit Status CANC wird in der Prognose so behandelt, als
        fände es noch statt. Das ist Absicht: Bei api-football ist CANC
        nicht verlässlich endgültig. Todesfelde gegen Norderstedt
        (Regionalliga Nord 2024/25) wurde am 21.04.2025 nachgeholt, steht
        in der API aber bis heute als CANC, ohne Ersatz-Spiel. Die
        ELO-Verlaufsseite zählt CANC nur bei der Frage, ob die Saison
        vorbei ist, als erledigt (Issue~\#292).
  \item[Rückzug eines Vereins]
        Zieht sich ein Verein während der Saison zurück (Türkgücü München,
        3.~Liga 2021/22; Türkspor Dortmund und KFC Uerdingen, Regionalliga
        West 2024/25), setzt die API seine Restspiele auf CANC. Seine
        bereits gespielten Spiele bleiben aber als beendet stehen, obwohl
        der Verband sie annulliert. Auch ein kampflos gewertetes Spiel kann
        als CANC statt als AWD kommen. Die Prognose bildet einen Rückzug
        deshalb nicht richtig ab; dafür braucht es einen Eingriff von Hand
        (Issue~\#292).
  \item[Kein Rückgang zur Mitte]
        Die ELO wird beim Saisonwechsel 1:1 übertragen. Viele ELO-Systeme
        ziehen die Werte zwischen zwei Saisons ein Stück zum Mittel zurück,
        um Kaderwechsel abzubilden.
  \item[Zwei Erwartungsskalen]
        ELO-Update und Tormodell rechnen mit unterschiedlich skalierten
        Erwartungen (Abschnitt~\ref{sec:skalen}).
  \item[Gewertete Spiele in der Playoff-ELO]
        Die aktuelle ELO für die Aufstiegsspiele
        (\code{rl\_aktuelle\_elo}, \datei{RCode/rl_verdrahtung.R:113--121})
        wird ohne das Kennzeichen \code{elo\_neutral} berechnet. Ein am grünen
        Tisch gewertetes Spiel würde dort die ELO bewegen, anders als
        überall sonst. Ob das beabsichtigt ist, ist im Code nicht kommentiert.
\end{description}


\part{Die Website}
\chapter{Seiten und Navigation}\label{kap:seiten}

\section{Grundprinzip: statische Seiten}

Die Website besteht aus statischen, in sich geschlossenen HTML-Seiten. Der
Scheduler schreibt sie nach jedem Simulationszyklus neu, ein Webserver
(Caddy) liefert sie aus. Es gibt keinen Anwendungsserver und keine Datenbank.
Bis August 2026 lief die Seite als Shiny-App auf shinyapps.io. ADR~0001
begründet den Wechsel: Der Dienst endet im März 2027, die App hatte keine
echte Interaktivität, die Server-Rechenzeit gebraucht hätte, und an einem
Spieltag waren bis zu 240 Neuveröffentlichungen nötig.

Die Heatmap ist eine HTML-Tabelle mit eingefärbten Zellen, es gibt keine
Bilder. Alle Seiten zusammen sind einige hundert Kilobyte groß. Nur das
ELO-Verlaufsdiagramm zeichnet ein Skript im Browser.

\section{Seiteninventar}

Pro Zyklus entstehen 22 Seiten
(\datei{RCode/generate_static_site.R:1086--1216}): zehn Ligaseiten, zehn
ELO-Verlaufsseiten, die Aufstiegsseite der Regionalligen und die
Methodik-Seite. Liegen gar keine Ergebnisse vor, entsteht nur eine
Platzhalter-Startseite (\enquote{Noch keine Prognosedaten verfügbar}).

\begin{center}
\footnotesize
\setlength{\tabcolsep}{4pt}
\begin{tabular}{@{}lllr@{}}
\toprule
Liga & Ligaseite & Verlaufsseite & ID \\
\midrule
Bundesliga & \code{index.html} & \code{bundesliga-verlauf.html} & 78 \\
2.~Bundesliga & \code{2-bundesliga.html} & \code{2-bundesliga-verlauf.html} & 79 \\
3.~Liga & \code{3-liga.html} & \code{3-liga-verlauf.html} & 80 \\
RL Nord & \code{rl-nord.html} & \code{rl-nord-verlauf.html} & 84 \\
RL Nordost & \code{rl-nordost.html} & \code{rl-nordost-verlauf.html} & 85 \\
RL West & \code{rl-west.html} & \code{rl-west-verlauf.html} & 87 \\
RL SüdWest & \code{rl-suedwest.html} & \code{rl-suedwest-verlauf.html} & 86 \\
RL Bayern & \code{rl-bayern.html} & \code{rl-bayern-verlauf.html} & 83 \\
Frauen-BL & \code{frauen-bundesliga.html} & \code{frauen-bundesliga-verlauf.html} & 82 \\
2.~Frauen-BL & \code{2-frauen-bundesliga.html} & \code{2-frauen-bundesliga-verlauf.html} & 1034 \\
\midrule
Aufstieg & \code{rl-aufstieg.html} & -- & -- \\
Methodik & \code{methodik.html} & -- & -- \\
\bottomrule
\end{tabular}
\end{center}

Die Bundesliga ist die Startseite und heißt deshalb \code{index.html}. Ihre
Verlaufsseite heißt \code{bundesliga-verlauf.html} und nicht
\code{index-verlauf.html}. Alle Namen stammen aus der Liga-Registry
(\datei{RCode/league_registry.R}) und aus \code{league\_views()}
(\datei{RCode/league_views.R}). Die Registry hält alles Ligaspezifische als
Daten (ADR~0005). Eine neue Liga ist dort ein neuer Eintrag und keine neue
Code-Verzweigung.

\section{Zweistufige Navigation}

Jede Seite trägt denselben Kopf: die Wortmarke \enquote{30 Punkte}, darunter
einen kurzen roten Strich, die Unterzeile und die Navigation. Die Navigation
hat eine Zeile je Gruppe. Jede Zeile beginnt mit dem Gruppennamen in
Kapitälchen, danach folgen die Links, getrennt durch Mittelpunkte:

\begin{center}
\begin{tabular}{@{}>{\sffamily\scriptsize\color{ink3}}l>{\sffamily\small}l@{}}
HERREN       & Bundesliga · 2.~Bundesliga · 3.~Liga \\
REGIONALLIGA & Nord · Nordost · West · SüdWest · Bayern · Aufstieg \\
FRAUEN       & Bundesliga · 2.~Bundesliga \\
             & Methodik \\
\end{tabular}
\end{center}

Die Regionalligen stehen bewusst direkt unter den Herren
(\datei{generate_static_site.R:54--69}, Issue~\#178). Sie sind
Herren-Ligen derselben Wechselgemeinschaft, und die Nachbarschaft soll das
zeigen. Die aktuelle Seite ist fett und unterstrichen und trägt
\code{aria-current="page"}.

Auf den Verlaufsseiten führen die Liga-Links zur \emph{Verlaufsseite} der
jeweils anderen Liga. Wer Verläufe vergleicht, bleibt so in dieser Ansicht.
Aufstieg und Methodik verlinken normal.

\section{Hinweis auf veraltete Daten}

Eine statische Seite veraltet genau dann, wenn der Scheduler nicht mehr
schreibt. Eine Prüfung beim Rendern würde deshalb nie anschlagen. Die
Entscheidung fällt im Browser (\datei{generate_static_site.R:445--469}):
Im Fuß jeder Seite steht der Zeitpunkt der Erzeugung als
\code{<time id="generated" datetime="…">} in UTC. Ein kleines Skript
berechnet daraus das Alter in Stunden,
\[
  \text{Alter} = \frac{\text{jetzt} - \text{erzeugt}}{3{,}6\cdot 10^{6}\ \text{ms}},
\]
und blendet ab einem Alter von mehr als 24 Stunden einen Hinweis ein:
\enquote{Achtung: Diese Prognosen sind \emph{n} Stunden alt und werden derzeit
nicht aktualisiert.}

Daneben gibt es einen zweiten, serverseitigen Hinweis. Arbeitet der
Scheduler wegen knappen API-Kontingents in einem gedrosselten Takt (mehr als
180 Sekunden statt 120), ergänzt der Fuß \enquote{Eingeschränkter Service,
Update etwa alle \emph{n} Minuten}. Anlass war ein Vorfall am 14.~September
2026, bei dem die Seite nur alle 83 Minuten aktualisiert wurde.

\chapter{Die Ligaseite, von oben nach unten}\label{kap:ligaseite}

Eine Ligaseite (\code{render\_league\_page},
\datei{RCode/generate_static_site.R:516--641}) besteht aus diesen Teilen:

\begin{enumerate}
  \item Kopf mit Wortmarke und Navigation (Kapitel~\ref{kap:seiten}),
  \item Hinweis auf veraltete Daten (verborgen, solange die Daten frisch sind),
  \item \textbf{Prognose}: Heatmap und zwei Zusammenfassungstabellen,
  \item \textbf{Tabelle}: Ligatabelle mit ELO, bei Regionalligen mit Zonen,
        darunter der Link zum ELO-Verlauf,
  \item \textbf{Rückblick}: die Spiele des aktuellen Spieltags mit Prognose und
        ELO-Änderung,
  \item \textbf{Live}: laufende Spiele, nur wenn welche laufen,
  \item \textbf{Ausblick}: die kommenden Spiele mit Wahrscheinlichkeiten und
        Ergebnis-Matrix,
  \item Fuß mit Zeitstempel und Blog-Link.
\end{enumerate}

\screenshot{ligaseite-bundesliga}{Kopf, Heatmap und Zusammenfassungstabellen der Bundesliga-Seite.}

\section{Prognose: die Heatmap}

Die Sektion beginnt mit \enquote{Saisonprognose \emph{Liga}} und dem Satz
\enquote{10\,000 Monte-Carlo-Simulationen des Saisonrests. Jede Zeile zeigt,
mit welcher Wahrscheinlichkeit ein Team auf welchem Platz landet.} Die
Heatmap ist die Matrix $\hat P$ aus Gleichung~\eqref{eq:mc}. Jede Zeile ist
eine Mannschaft, sortiert nach erwartetem Platz, jede Spalte ein Endplatz.
Die Zeilenköpfe zeigen das Kürzel. Der volle Vereinsname erscheint als
Tooltip, beim Antippen auch auf Touch-Geräten.

\subsection*{Die Farbformel}

Jede Zelle wird von Weiß zum Tintenblau \code{\#1F3A4D} interpoliert
(\code{.heat\_style}, \datei{RCode/render_helpers.R:83--96}). Damit kleine
Wahrscheinlichkeiten sichtbar bleiben und nicht linear verblassen, wird die
Wahrscheinlichkeit vorher mit dem Exponenten 0{,}6 gestaucht:
\begin{equation}
  t = p^{\,0{,}6}, \qquad
  \begin{pmatrix} R\\ G\\ B \end{pmatrix}
  = \operatorname{round}\!\left(
    \begin{pmatrix} 255\\ 255\\ 255 \end{pmatrix}
    + t\left[\begin{pmatrix} 31\\ 58\\ 77 \end{pmatrix} - \begin{pmatrix} 255\\ 255\\ 255 \end{pmatrix}\right]
  \right).
\end{equation}
Ist $t > 0{,}52$, also ab $p \approx 33{,}6\,\%$, wechselt die Schrift von
fast Schwarz (\code{\#15130F}) zu Weiß.

\begin{center}
\definecolor{h01}{RGB}{241,243,244}
\definecolor{h10}{RGB}{199,206,210}
\definecolor{h50}{RGB}{107,125,138}
\definecolor{h100}{RGB}{31,58,77}
\begin{tabular}{@{}lcccc@{}}
\toprule
$p$ & 1\,\% & 10\,\% & 50\,\% & 100\,\% \\
$t = p^{0{,}6}$ & 0{,}063 & 0{,}251 & 0{,}660 & 1 \\
RGB & (241,\,243,\,244) & (199,\,206,\,210) & (107,\,125,\,138) & (31,\,58,\,77) \\
Farbe &
\colorbox{h01}{\makebox[2.6em]{\ttfamily\small 1}} &
\colorbox{h10}{\makebox[2.6em]{\ttfamily\small 10}} &
\colorbox{h50}{\makebox[2.6em]{\ttfamily\small\color{white} 50}} &
\colorbox{h100}{\makebox[2.6em]{\color{white}\checkmark}} \\
\bottomrule
\end{tabular}
\end{center}

\subsection*{Was in einer Zelle steht}

Eine Zelle mit $p = 0$ bleibt leer und weiß. Sonst steht dort der
Prozentwert nach der Regel aus Abschnitt~\ref{sec:prozent}, ohne
Prozentzeichen, denn die Legende sagt \enquote{Angaben in Prozent}. Auf
schmalen Bildschirmen (bis 520~px) werden Werte unter einem Prozent nur noch
als Färbung gezeigt. Die Heatmap passt so ohne seitliches Scrollen auf ein
Telefon.

\section{Prognose: die Zusammenfassungstabellen}

Unter der Heatmap stehen zwei Tabellen nebeneinander, oben die Spitze und
unten das Tabellenende. Jede Spalte summiert eine Gruppe von Plätzen.
Gezeigt werden nur Mannschaften, deren Summe über die Spalten mindestens
1~Prozent beträgt. Eine 0 erscheint hier als \enquote{0} und nicht als
leere Zelle wie in der Heatmap.

\begin{table}[ht]
\centering\small
\caption{Spaltengruppen je Liga (\texttt{RCode/league\_views.R}). Negative
Zahlen zählen von unten. \enquote{berechnet} heißt, dass der Wert nicht aus
einer Platzsumme stammt, sondern aus Kapitel~\ref{kap:regionalliga}.}
\label{tab:panels}
\begin{tabularx}{\textwidth}{@{}l>{\raggedright}X>{\raggedright\arraybackslash}X@{}}
\toprule
Liga & oben & unten \\
\midrule
Bundesliga & Meister (1), Champions League (2--4), Europa League (5), Conference League Quali (6) & Relegation (16), Abstieg (17--18) \\
2.~Bundesliga & Aufstieg (1--2), Relegation Bundesliga (3) & Relegation 3.~Liga (16), Abstieg (17--18) \\
3.~Liga & aus dem Malus-Lauf: Aufstieg (1--2), Relegation (3), DFB-Pokal (4) & Abstieg (17--20) \\
Frauen-Bundesliga & Meister (1), Champions League Quali (2--3) & Abstieg ($-2$ bis $-1$) \\
2.~Frauen-Bundesliga & aus dem Malus-Lauf: Aufstieg (1--2) & Abstieg ($-3$ bis $-1$) \\
RL Nord, RL Bayern & Meister (1), Aufstieg (berechnet, Aufstiegsspiele) & Abstieg (berechnet); Bayern zusätzlich Relegation \\
RL Nordost, West, SüdWest & Aufstieg (1, Direktaufstieg) & Abstieg (berechnet) \\
\bottomrule
\end{tabularx}
\end{table}

Berechnete Spalten werden über ihren Namen zugeordnet, nicht über ihre
Position. Fehlt eine, bricht das Rendern dieser Seite ab, statt eine falsche
Zahl zu zeigen (\datei{generate_static_site.R:246--285}).

\section{Tabelle: Ligatabelle und ELO}

Die Sektion \enquote{Ligatabelle und ELO} zeigt den aktuellen Stand, daneben
die Stärkeschätzung des Modells und ihre Veränderung seit Saisonbeginn
(\code{render\_liga\_tabelle}, \datei{RCode/render_sections.R:115--164}).

\begin{center}
\small
\begin{tabular}{@{}lll@{}}
\toprule
Spalte & Inhalt & sortierbar \\
\midrule
Platz    & Tabellenplatz & ja \\
Team     & voller Name, Kürzel als Tooltip & -- \\
Sp.      & Spiele (auf Telefonen ausgeblendet) & -- \\
Tordiff. & Tordifferenz mit Vorzeichen (auf Telefonen ausgeblendet) & -- \\
Punkte   & Punkte & ja \\
ELO      & aktuelle ELO, eine Nachkommastelle & ja \\
$\Delta$ ELO & ELO heute minus ELO zum Saisonstart & ja \\
\bottomrule
\end{tabular}
\end{center}

Die Tabelle zählt beendete und gewertete Spiele und sortiert nach Punkten,
Tordifferenz und Toren. Ein Klick auf einen Spaltenkopf sortiert im Browser
um. Die Sortierung nach $\Delta$ ELO zeigt oben, wer die Erwartung vom
Saisonstart am deutlichsten übertrifft (Issue~\#285). Vorzeichen werden vom
\emph{gerundeten} Wert bestimmt: \enquote{$+$}, \enquote{$-$} (echtes
Minuszeichen U+2212) oder \enquote{$\pm$} bei null.

\subsection*{Auf- und Abstiegszonen der Regionalligen}

In den Regionalligen markiert eine dünne farbige Linie am linken Zeilenrand
die Zonen: \textcolor{akzent}{Rot} für Abstieg, \textcolor{ocker}{Ocker} für
Relegation (nur Bayern), \textcolor{gruen}{Grün} für Aufstieg. Die
\emph{Deckkraft} der Linie ist die Wahrscheinlichkeit selbst, mindestens aber
15~Prozent, damit jede Chance über null sichtbar bleibt. Treffen mehrere
Zonen zusammen, gilt der Vorrang Rot vor Ocker vor Grün. Die Zonen gehören
zum Tabellenplatz und verschwinden, sobald nach einer anderen Spalte sortiert
wird.

Darunter erklärt eine Fußnote die Regel der Staffel, etwa für Nord:
\enquote{Drei Vereine steigen ab, je Absteiger aus der 3.~Liga einer mehr.
Gewinnt Nord das Aufstiegsspiel, ist es einer weniger.} Schwankt die Zahl
der Absteiger, folgt ihre Verteilung (\enquote{3 mit $x$\,\%, 4 mit
$y$\,\% …}) und danach die Wahrscheinlichkeit je Platz.

Unter der Tabelle führt ein Link zur ELO-Verlaufsseite der Liga
(Abschnitt~\ref{sec:verlauf}).

\section{Rückblick}

Der Rückblick zeigt die beendeten Spiele des aktuellen Spieltags. Solange
dieser nicht abgeschlossen ist, zeigt er zusätzlich den vorigen und dazu
gekennzeichnete Nachholspiele. Jede Zeile hat fünf Spalten:

\begin{enumerate}
  \item Anstoß in Berliner Zeit, etwa \enquote{Sa. 4.10., 15:30 Uhr},
  \item die Paarung \enquote{Heim – Gast},
  \item einen 1/X/2-Balken mit den Wahrscheinlichkeiten \emph{vor} dem Spiel:
        \textcolor{tinte}{Heimsieg}, \textcolor{ink3}{Remis} (hellgrau),
        \textcolor{ocker}{Auswärtssieg},
  \item das Ergebnis,
  \item die ELO-Änderung als \enquote{Heim / Gast}, etwa
        \enquote{$+4{,}1$ / $-4{,}1$}.
\end{enumerate}

Die Wahrscheinlichkeiten des Balkens stammen aus Abschnitt~\ref{sec:tor-beispiel}
und werden mit der ELO unmittelbar vor dem Spiel berechnet. Gerundet wird so,
dass die drei Segmente zusammen genau 100 ergeben:
\begin{equation}
  h = \operatorname{round}(100\,p_H), \qquad a = \operatorname{round}(100\,p_A), \qquad x = 100 - h - a.
\end{equation}
Segmente unter 8~Prozent tragen keine Zahl. Gewertete Spiele sind mit
\enquote{Wertung} markiert und haben keine ELO-Änderung.

\section{Live und Ausblick}

\textbf{Live} erscheint nur, wenn Spiele laufen. Die Sektion zeigt Paarung
und Zwischenstand mit dem Hinweis \enquote{Prognosen werden während des
Spiels nicht aktualisiert.}

\textbf{Ausblick} zeigt die offenen Spiele des aktuellen und des nächsten
Spieltags, gerechnet mit der ELO von heute, jeweils mit 1/X/2-Balken. Zu
jeder Paarung lässt sich die \emph{Ergebnis-Matrix} aufklappen
(\datei{render_sections.R:561--587}). Sie zeigt das $7\times7$-Gitter aus
Kapitel~\ref{kap:tormodell}: Zeilen sind die Tore der Heimmannschaft, Spalten
die der Gastmannschaft, und die letzte Zeile und Spalte fassen \enquote{6+}
zusammen. Die Werte haben eine Nachkommastelle, Zellen unter 0{,}1~Prozent
bleiben leer. Die Färbung ist relativ zum größten Wert der Matrix:
\code{.heat\_style}$\bigl(\min(1,\ 0{,}75 \cdot p / p_{\max})\bigr)$. Auch
das wahrscheinlichste Ergebnis wird also nie ganz dunkel.

\section{Fuß}

Links steht \enquote{Letztes Update: TT.MM.JJJJ HH:MM MESZ} in Berliner Zeit,
dahinter der maschinenlesbare Zeitstempel für den Veraltet-Hinweis, bei
gedrosseltem Takt zusätzlich \enquote{Eingeschränkter Service}. Rechts steht
\enquote{Mehr dazu unter 30punkte.wordpress.com}.

\chapter{Aufstieg, Methodik, ELO-Verlauf}\label{kap:weitere}

\section{Die Aufstiegsseite der Regionalligen}

Die Seite \code{rl-aufstieg.html} gehört zu keiner Liga. Sie fasst den
Aufstieg aus allen fünf Staffeln zusammen
(\code{.render\_aufstiegsseite}, \datei{RCode/generate_static_site.R:860--904}).
Unter der Überschrift \enquote{Aufstieg in die 3.~Liga} erklärt ein kurzer
Text die Regel. Danach folgt je Staffel eine Tabelle mit drei Spalten:

\begin{description}[font=\sffamily\bfseries\small]
  \item[Meister] $P(\text{Platz 1})$ aus dem Hauptlauf der Staffel.
  \item[Aufstieg] Bei Direktaufsteigern gleich dem Meister, bei den
        Playoff-Staffeln nach Gleichung~\eqref{eq:rl-aufstieg}.
  \item[Siegquote] $P(\text{Aufstieg})/P(\text{Meister})$, nur bei den
        Playoff-Staffeln. Leer, wenn der Quotient nicht definiert ist.
\end{description}

Welche Staffel Playoff-Staffel ist, liest die Seite aus der Ansicht ab und
nicht aus festen Staffelnamen. So folgt sie der jährlichen Rotation ohne
Codeänderung. Gezeigt werden nur Mannschaften mit einer Meister- oder
Aufstiegschance über null.

\screenshot{aufstiegsseite}{Die Aufstiegsseite mit den fünf Staffeln.}

\section{Die Methodik-Seite}

Der Text der Methodik-Seite ist redaktionell und liegt als fertiges HTML in
\datei{RCode/site_assets/methodik_content.html}. Er geht auf den Blogbeitrag
\enquote{Was die Prognosen mit Schach zu tun haben} vom 1.~April 2015 zurück
und wurde für den Relaunch überarbeitet. Er erklärt das Modell ohne Formeln:

\begin{enumerate}
  \item Die Stärkeschätzung ELO, Heimvorteil 40 Punkte, gedämpfte Wirkung
        hoher Siege,
  \item vom Stärkeunterschied zum Ergebnis: Torerwartung, Poisson-Verteilung,
        das Beispiel Bayern gegen Stuttgart,
  \item zehntausend Saisons: das Monte-Carlo-Verfahren,
  \item Abstieg aus den Regionalligen und warum er erst am Saisonende
        feststeht,
  \item was die Prognose kann und was nicht: Häkchen, Gleichstand,
        Zweitvertretungen, Wertungen, Modellgüte.
\end{enumerate}

Ein farbig abgesetzter Absatz vermerkt die Neukalibrierung des
Heimvorteils von 65 auf 40 Punkte im September 2026.

\section{Die ELO-Verlaufsseiten}\label{sec:verlauf}

Seit Oktober 2026 hat jede Liga eine Verlaufsseite (Issue~\#184). Sie zeigt,
wie sich die Stärkeschätzung jeder Mannschaft über die Saison entwickelt hat.
Gestaltet ist sie nach den Grundsätzen von Edward Tufte, mit den
Gestaltungsmitteln der Website.

\screenshot{verlauf-bundesliga}{ELO-Verlauf mit hervorgehobenem Verein und Tooltip.}

\subsection*{Aufbau}

\begin{enumerate}
  \item Kopf und Navigation (die Liga-Links führen zu den anderen
        Verlaufsseiten), darunter ein Rücklink \enquote{← Prognose \emph{Liga}}.
  \item Überschrift \enquote{ELO-Verlauf \emph{Liga} \emph{Saison}} und der
        Einleitungssatz \enquote{Jede Linie ist ein Verein, jeder Knick ein
        Spiel. […] Die Zahl am Ende jeder Linie ist der heutige Stand.}
  \item Das Diagramm mit einem Umschalter \enquote{Nach Spielen} /
        \enquote{Nach Datum}.
  \item Ein Hinweis, dass die Kurven aus denselben Spielen und demselben
        Modell entstehen wie die Prognose, mit Verweis auf die Methodik. Er
        nennt bewusst keine Modellkonstanten.
  \item Die Tabelle \enquote{Stärkeschätzung heute}: Verein, ELO zum
        Saisonbeginn, ELO heute und Veränderung, sortiert nach der heutigen
        ELO. Sie wird auf dem Server erzeugt und funktioniert auch ohne
        JavaScript.
\end{enumerate}

\subsection*{Die Daten}

\code{elo\_verlauf\_daten()} (\datei{RCode/elo_verlauf.R:41--116}) baut für
jeden Verein eine Punktfolge. Punkt 0 ist die ELO zum Saisonstart. Danach
folgt je \emph{gespieltem} Spiel ein Punkt mit Datum, Spieltag,
Nachholspiel-Kennzeichen, Gegner, Ergebnis, ELO vor und nach dem Spiel, der
Änderung $\delta$ und den Wahrscheinlichkeiten für Sieg und Remis vor dem
Spiel. Gewertete Spiele erzeugen keinen Punkt, weil sie die ELO nicht
bewegen.

Eine Invariante sichert die Konsistenz mit der Ligatabelle: Der letzte Punkt
jeder Linie muss auf $10^{-6}$ genau der ELO in der Tabelle entsprechen.
Sonst entfällt diese eine Verlaufsseite mit einer Warnung im Log, und alle
anderen Seiten entstehen trotzdem.

Die Daten gehen als JSON in einem
\code{<script type="application/json">}-Element in die Seite. Jedes
\enquote{\textless} wird dabei als \code{\textbackslash u003c} kodiert. Das
bleibt gültiges JSON, und kein Vereinsname kann das Skript-Element mit
\code{</script>} oder \code{<!--} vorzeitig beenden.

\subsection*{Das Diagramm}

Das Skript \datei{RCode/site_assets/verlauf.js} zeichnet ein SVG mit dem
Koordinatensystem $1080 \times 640$.

\begin{description}[font=\sffamily\bfseries\small]
  \item[Wenig Tinte] Keine Rahmen, Achsenlinien, Schatten oder Flächen.
        Waagerechte Hilfslinien liegen alle 50 ELO-Punkte. Alle Linien sind
        leise grau, genau eine ist hervorgehoben (Tintenblau, 2~px). Es gibt
        keine Vereinsfarben und keine Legende. Jede Linie ist an ihrem Ende
        direkt beschriftet, mit dem Kürzel und der gerundeten heutigen ELO.
  \item[Senkrechte Skala] Sie reicht vom kleinsten bis zum größten Wert, mit
        einem Polster von 6~Prozent der Spanne, mindestens 10 Punkte. Liegt
        die Oberkante nah an einer 50er-Linie, rastet sie dort ein, damit das
        Diagramm oben mit einer Hilfslinie abschließt.
  \item[Zwei Zeitachsen] \enquote{Nach Spielen} (Standard) setzt das
        $n$-te Spiel jedes Vereins an dieselbe Stelle. \enquote{Nach Datum}
        zeigt die echte Zeitachse mit Monatsnamen und markiert die
        Winterpause als Fläche, wenn die größte Lücke länger als 18 Tage ist.
        Bei laufender Saison bleibt rechts Platz für den nächsten Spieltag.
  \item[Beschriftungen] Die Endbeschriftungen stehen in einer gemeinsamen
        Spalte rechts und sind auf mindestens 12{,}5~px Abstand
        auseinandergezogen. Vereine mit weniger Spielen erreichen die Spalte
        über eine gepunktete Führungslinie.
  \item[Tooltip] Zeigen oder Tippen auf eine Linie wählt den Verein und das
        letzte Spiel links vom Zeiger. Der Tooltip zeigt Datum und Spieltag,
        Paarung und Ergebnis, \enquote{Das Modell gab \emph{X} vorher
        $a$\,\% Sieg, $b$\,\% Remis.} und \enquote{ELO vor → nach $\pm\delta$}.
  \item[Tabelle und Diagramm] Fährt die Maus über eine Tabellenzeile, wird
        die zugehörige Linie hervorgehoben. Ein Klick hält die Auswahl.
        Umgekehrt markiert die Auswahl im Diagramm die Tabellenzeile.
  \item[Telefon] Das Diagramm bleibt mindestens 640~px breit und lässt sich
        seitlich scrollen (\code{touch-action: pan-x pan-y}).
\end{description}

\chapter{Gestaltung und Formatregeln}\label{kap:design}

Die Gestaltung ist zurückhaltend und zeitungsnah. Die Daten haben Vorrang,
und Farbe trägt Bedeutung, statt zu schmücken. Alle Regeln stehen in einer
handgeschriebenen Stilvorlage, \datei{RCode/site_assets/site.css}. Sie
kommt ohne Framework aus.

\section{Schriften}

Alle drei Schriften liegen als WOFF2 im Projekt, nur mit lateinischem
Zeichensatz, und werden selbst ausgeliefert. Es gibt keine Anfragen an
Schriftendienste.

\begin{center}
\small
\begin{tabularx}{\textwidth}{@{}llX@{}}
\toprule
Rolle & Schrift & Einsatz \\
\midrule
Serif & Source Serif 4 (variabel, 200--900) & Fließtext, Überschriften, Wortmarke, Vereinsnamen; Mediävalziffern im Text \\
Sans  & Inter Tight (400--700) & Navigation, Kapitälchen-Überschriften, Tabellenköpfe, Legenden, Fuß \\
Mono  & JetBrains Mono (400--700) & alle Zahlen in Tabellen, Heatmap, Ergebnisse, ELO (Tabellenziffern) \\
\bottomrule
\end{tabularx}
\end{center}

Dieses Handbuch ist in Source Serif Pro, Inter und Source Code Pro gesetzt,
den nächsten Verwandten in \TeX\ Live.

\section{Farben}

\newcommand{\farbe}[3]{\textcolor{#1}{\rule{1.6em}{0.9em}} & \code{#2} & \code{#3}}
\begin{center}
\small
\begin{tabularx}{\textwidth}{@{}cllX@{}}
\toprule
& Token & Wert & Verwendung \\
\midrule
\farbe{ink}{--ink}{\#15130F} & Haupttext \\
\farbe{ink2}{--ink-2}{\#312E29} & Sekundärtext, Navigation \\
\farbe{ink3}{--ink-3}{\#6B655C} & Legenden, Fuß, Kapitälchen \\
\farbe{rule}{--rule}{\#D4CEBF} & Linien, Rahmen, Linien im Verlauf \\
\farbe{rulesoft}{--rule-soft}{\#E3DED0} & feine Linien, Remis-Segment \\
\farbe{surface2}{--surface-2}{\#F5F3EE} & Hinweisflächen, Winterpause, gewählte Zeile \\
\farbe{akzent}{--akzent}{\#B0261E} & Rot: Unterstreichung, Strich unter der Wortmarke, Abstieg, negative Änderung, Fokusrahmen \\
\farbe{tinte}{--tinte}{\#1F3A4D} & Tintenblau: Endfarbe der Heatmap, Heimsieg, hervorgehobene Linie, positive Änderung \\
\farbe{ocker}{--ocker}{\#B8893A} & Auswärtssieg, Relegation, Hinweisbalken \\
\farbe{gruen}{--gruen}{\#3F7D3A} & Aufstiegszone \\
\bottomrule
\end{tabularx}
\end{center}

Die Seite hat bewusst nur ein helles Erscheinungsbild. Einen Dunkelmodus
gibt es nicht.

\section{Satz und Raster}

\begin{itemize}
  \item Fließtext: Serif, 17~px, Zeilenhöhe 1{,}58. Der Inhalt ist höchstens
        1\,080~px breit, Einleitungstexte höchstens 64 Zeichen.
  \item Abschnitte: 56~px Abstand und eine feine Linie. Jede Sektion beginnt
        mit einer kleinen Überschrift in Kapitälchen (\enquote{PROGNOSE},
        \enquote{TABELLE}, \dots), danach folgt die Überschrift in Serif.
  \item Links: dunkle Schrift mit roter Unterstreichung.
  \item Umbrüche bei \textbf{760~px} (kleinere Abstände, Spielzeilen
        zweispaltig, Verlaufsdiagramm scrollbar) und \textbf{520~px}
        (Ligatabelle ohne \enquote{Sp.} und \enquote{Tordiff.}, Heatmap mit
        fester Spaltenbreite, Werte unter 1~Prozent nur als Farbe).
\end{itemize}

\section{Zahlen und Prozente}\label{sec:prozent}

Prozentwerte formatiert eine einzige Funktion, \code{prozent()}
(\datei{RCode/render_helpers.R:32--47}):

\begin{center}
\small
\begin{tabular}{@{}ll@{}}
\toprule
Wahrscheinlichkeit $x$ & Anzeige \\
\midrule
$x = 0$ & \enquote{0} (in der Heatmap: leere Zelle) \\
$0 < x < 0{,}01$ & \enquote{<1} \\
$0{,}01 \le x \le 0{,}99$ & $\operatorname{round}(100x)$ \\
$0{,}99 < x < 1$ & \enquote{>99} \\
$x = 1$ & \enquote{\checkmark} \\
\bottomrule
\end{tabular}
\end{center}

Damit steht nie \enquote{100} oder \enquote{0} für etwas, das in einzelnen
Simulationen doch anders ausging. Weitere Regeln:

\begin{itemize}
  \item Dezimalkomma überall. Negative Zahlen tragen das echte Minuszeichen
        (U+2212), null trägt \enquote{$\pm$}.
  \item Wo ein Prozentzeichen steht, trennt es ein geschütztes Leerzeichen von
        der Zahl (\enquote{12\,\%}). Tabellen verzichten auf das Zeichen und
        sagen es in der Legende.
  \item Die Ergebnis-Matrix zeigt eine Nachkommastelle, alle anderen
        Prozentangaben sind ganzzahlig.
  \item Zeiten stehen in Berliner Zeit mit \enquote{Uhr}, im Fuß mit
        \enquote{MEZ} oder \enquote{MESZ}.
\end{itemize}

\section{Barrierefreiheit}

\begin{itemize}
  \item Die Heatmap trägt \code{aria-label}, und die Balken sind mit
        \code{role="img"} und einer gesprochenen Beschreibung versehen
        (\enquote{Sieg Heim 62\,\%, Remis 22\,\%, Sieg Gast 16\,\%}).
  \item Kürzel sind als \code{<abbr>} mit dem vollen Namen ausgezeichnet und
        per Tastatur und Antippen erreichbar.
  \item Das Verlaufsdiagramm hat eine Textbeschreibung, und die Tabelle
        darunter enthält dieselben Zahlen als zugängliche Alternative.
  \item Alle bedienbaren Elemente haben einen sichtbaren roten Fokusrahmen.
        Die einzige Animation (das Aufklappen der Ergebnis-Matrix) entfällt
        bei \code{prefers-reduced-motion}.
\end{itemize}


\part{Technik}
\chapter{Architektur}\label{kap:architektur}

\section{Komponenten und Datenfluss}

Das System besteht aus einem einzigen Docker-Container, in dem zwei Prozesse
laufen: der R-Scheduler und der Rust-Simulationsserver. Dazu kommen ein
Docker-Volume und ein Webserver.

\begin{figure}[ht]
\centering
\resizebox{\textwidth}{!}{\begin{tikzpicture}[
  font=\sffamily\footnotesize,
  box/.style={draw=ink3,fill=white,rounded corners=1.5pt,align=center,inner sep=5pt,minimum height=1.05cm},
  ext/.style={box,fill=surface2},
  arr/.style={-{Stealth[length=5pt]},draw=ink2,thick},
  lab/.style={font=\sffamily\scriptsize,color=ink3,align=center},
  grp/.style={draw=rule,dashed,rounded corners=3pt,inner sep=8pt},
]
\node[ext,text width=2.1cm] (api) {API-Football\\[1pt]\scriptsize(RapidAPI)};
\node[box,right=1.3cm of api,text width=2.6cm] (sch) {\textbf{R-Scheduler}\\[1pt]\scriptsize 11--23 Uhr, Takt 2 min};
\node[box,below=1.15cm of sch,text width=2.6cm] (rust) {\textbf{Rust-Server}\\[1pt]\scriptsize axum + rayon, :8080};
\node[box,right=1.6cm of sch,text width=2.4cm] (gen) {\textbf{Seitengenerator}\\[1pt]\scriptsize R + htmltools};
\node[ext,right=1.5cm of gen,text width=2.1cm,shape=cylinder,shape border rotate=90,aspect=0.18,minimum height=1.25cm] (vol) {Volume\\[1pt]\scriptsize\mbox{\texttt{fussball-site}}};
\node[box,below=2.8cm of vol,text width=2.1cm] (caddy) {\textbf{Caddy}\\[1pt]\scriptsize HTTPS, Cache};
\node[ext,below=2.8cm of gen,text width=2.4cm] (browser) {Browser\\[1pt]\scriptsize\mbox{\texttt{verlauf.js}}, Veraltet-Hinweis};
\node[ext,below=1.15cm of api,text width=2.1cm] (team) {TeamList-CSV\\[1pt]\scriptsize Start-ELO};

\begin{scope}[on background layer]
  \node[grp,fit=(sch)(rust)(gen),label={[lab,anchor=south west]north west:Container \texttt{fussball-scheduler}}] {};
\end{scope}

\draw[arr] (api) -- node[lab,above]{Spiele} (sch);
\draw[arr] ([xshift=-6pt]sch.south) -- node[lab,left]{JSON} ([xshift=-6pt]rust.north);
\draw[arr] ([xshift=6pt]rust.north) -- node[lab,right]{Matrix,\\Details} ([xshift=6pt]sch.south);
\draw[arr] (sch) -- node[lab,above]{Ergebnisse} (gen);
\draw[arr] (gen) -- node[lab,above]{22 Seiten} (vol);
\draw[arr] (vol) -- node[lab,right]{nur lesend} (caddy);
\draw[arr] (caddy) -- node[lab,above]{HTTPS} (browser);
\draw[arr,dashed] (team) -- (sch);
\end{tikzpicture}}
\caption{Komponenten zur Laufzeit. Der Container ist gestrichelt umrandet.}
\label{fig:architektur}
\end{figure}

\begin{enumerate}
  \item \textbf{API-Football} über RapidAPI liefert Spielpläne, Ergebnisse
        und Live-Status.
  \item Der \textbf{R-Scheduler} fragt die API im Zwei-Minuten-Takt ab,
        bereitet die Daten auf und ruft den Rust-Server.
  \item Der \textbf{Rust-Server} rechnet die Monte-Carlo-Simulation und die
        deterministischen Spieldetails. Er ist nur innerhalb des Containers
        erreichbar.
  \item Der \textbf{Seitengenerator} (R, \code{htmltools}) schreibt die 22
        Seiten in das Ausgabeverzeichnis. Die Ergebnisse übergibt der
        Scheduler im Speicher. Im Produktionspfad wird keine Ergebnisdatei
        geschrieben.
  \item Das \textbf{Docker-Volume} \code{fussball-site} teilen Container und
        Webserver. Der Container schreibt, Caddy liest.
  \item \textbf{Caddy} liefert die Seiten mit HTTPS, Kompression und
        passenden Cache-Regeln aus.
\end{enumerate}

\section{Leitentscheidungen}

Die Architektur folgt einer Reihe festgehaltener Entscheidungen
(\datei{docs/adr/}):

\begin{center}
\small
\begin{tabularx}{\textwidth}{@{}lX@{}}
\toprule
ADR & Entscheidung \\
\midrule
0001 & Statische Seiten statt einer gehosteten Shiny-App. \\
0002 & Spieldetails werden bei jedem Rendern deterministisch aus Startwerten und Ergebnissen berechnet und nicht gespeichert. Die Modellkonstanten leben nur im Rust-Server. \\
0003 & ELO-Eichung für Ligen ohne Historie. \\
0004 & Ein Tormodell je Wechselgemeinschaft. \\
0005 & Ligakonfiguration als Daten in einer Registry. \\
0006 & Abstiegskopplung der Regionalligen an die 3.~Liga, exakt ausgezählt. \\
0007 & Die TeamList ist ein gepflegtes Stammdatenblatt, der Saisonwechsel läuft in zwei Phasen. \\
0008 & Zweitvertretungs-Malus nach §\,55b Nr.~3.1 DFB-Spielordnung. \\
\bottomrule
\end{tabularx}
\end{center}

Eine frühere Aufteilung in Microservices mit Kubernetes wurde im Mai 2026
zugunsten des einen Containers zurückgenommen (Git-Tag
\code{pre-deployment-cleanup-2026-05-02}). Diagramme dieses Entwurfs in
\datei{docs/architecture/} sind als historisch markiert.

\chapter{Scheduler und Datenbeschaffung}\label{kap:scheduler}

Der Scheduler ist ein R-Prozess (\datei{RCode/updateScheduler.R}, Hauptschleife
in \datei{RCode/update_all_leagues_loop.R}). Er läuft dauerhaft im
Container und entscheidet, wann Daten geholt, simuliert und Seiten
geschrieben werden.

\section{Das Zeitfenster}

\begin{itemize}
  \item Aktiv ist der Scheduler täglich von \textbf{11:00 bis 23:00 Uhr}
        Berliner Zeit, also 720 Minuten. Der Beginn um 11 Uhr folgt aus dem
        Spielplan: 37~Prozent der Spiele der 2.~Frauen-Bundesliga beginnen
        vormittags um 11 Uhr.
  \item Die Laufzeit ergibt sich aus Fensterende minus Fensterbeginn. Die
        Umgebungsvariable \code{DURATION} bleibt leer. Ein fester Wert wäre
        eine zweite Quelle, die beim nächsten Verschieben des Fensters
        unbemerkt das Fenster beschneiden würde.
  \item Außerhalb des Fensters schläft der Prozess bis zum nächsten
        Fensterbeginn. Jeder Prozess ist ein Tageslauf: Er endet, sobald die
        nächste Runde nicht mehr vor 23 Uhr passt, und Docker startet den
        Container neu (\code{restart: unless-stopped}). Der neue Prozess
        schläft dann bis 11 Uhr -- pro Abend genau ein Neustart.
  \item Als „nach dem Fenster“ gilt deshalb schon ein Start ab
        \textbf{22:58 Uhr}, also wenn höchstens noch ein Takt (zwei Minuten)
        bleibt. Vorher plante ein Start um 22:58 bis 23:00 noch eine Runde,
        endete nach ihr und startete erneut -- rund sieben Neustarts je Abend,
        die \code{RestartCount} als Absturzanzeige unbrauchbar machten
        (Issue~\#287).
  \item Die Fensterrechnung erfolgt ausdrücklich in \code{Europe/Berlin},
        unabhängig von der Zeitzone des Prozesses.
\end{itemize}

\section{Takt und Live-Gating}

Im Normalbetrieb läuft eine Runde alle \textbf{120 Sekunden}, das sind bis
zu 361 Runden am Tag. Damit das API-Kontingent reicht, ist jede Runde zweistufig:

\begin{enumerate}
  \item \textbf{Live-Abfrage} (eine Anfrage für alle Ligen,
        \code{GET /v3/fixtures?live=78-79-80-…}): Läuft irgendwo ein Spiel?
  \item \textbf{Vollabruf} (eine Anfrage je Liga, also zehn) nur, wenn
        \begin{itemize}
          \item ein Spiel läuft, oder
          \item ein Spiel den Live-Feed verlassen hat, in den Saisondaten aber
                noch nicht als beendet steht (Issue~\#154), oder
          \item der letzte erfolgreiche Vollabruf länger als 3\,600 Sekunden
                zurückliegt (Sicherheitsnetz, Issue~\#226), oder
          \item die Live-Abfrage scheitert.
        \end{itemize}
\end{enumerate}

Die erste Runde eines Tages ist immer ein Vollabruf. Simuliert wird eine
Liga nur, wenn sich die Menge ihrer beendeten Spiele geändert hat. Verglichen
wird die Menge der Spiel-IDs, nicht ihre Anzahl. Seiten werden geschrieben,
wenn simuliert wurde oder sich etwas Sichtbares geändert hat, etwa ein
Zwischenstand oder eine Anstoßzeit.

\section{Das API-Kontingent}

Der Zugang läuft über den Pro-Plan von API-Football: 7\,500 Anfragen am Tag
und 300 je Minute. Ein ruhiger Tag verbraucht etwa 500 Anfragen, ein voller
Spieltag bis zu 3\,900 (Issue~\#190).

Nach jeder Anfrage liest der Scheduler die Kopfzeilen
\code{x-ratelimit-requests-remaining}, \code{-limit} und \code{-reset} mit.
Ein Regler (\datei{RCode/rate_limit_takt.R}) passt den Takt an das
Restkontingent an. Er hält den Abstand zwischen 60 und 5\,400 Sekunden,
plant mit einem Sicherheitsabschlag von 10~Prozent und schaltet mit 15~Prozent
Hysterese, damit der Takt nicht flattert. Bei gedrosseltem Takt zeigt die
Website einen Hinweis (Kapitel~\ref{kap:seiten}).

\section{Fehlertoleranz}

\begin{itemize}
  \item \textbf{Je Liga isoliert.} Scheitert der Abruf, die Aufbereitung oder
        die Simulation einer Liga, fällt nur diese Liga aus. Ihre letzte
        Prognose bleibt stehen, und alle anderen Seiten werden aktualisiert.
  \item \textbf{Zeitlimits.} 30 Sekunden für API-Football, 60 Sekunden für den
        Rust-Server, gegen hängende Verbindungen.
  \item \textbf{Startskript.} \datei{docker-start.sh} startet den Rust-Server,
        wartet auf dessen \code{/health}, startet ihn neu, falls er stirbt,
        und versucht den Scheduler bis zu fünfmal im Abstand von 30 Sekunden
        zu starten.
  \item \textbf{Warnungen sofort.} \code{options(warn = 1)} schreibt jede
        Warnung direkt ins Container-Log. Viele Fehler führen nicht zum
        Absturz, sondern zu einer fehlenden Sektion mit Warnung.
\end{itemize}

\section{Der Ablauf einer Runde}

\begin{enumerate}
  \item Kontingent prüfen, Takt festlegen.
  \item Live-Abfrage, bei Bedarf Vollabruf je Liga.
  \item Daten aufbereiten (\code{transform\_data}): Statusgruppen, Filter auf
        Hauptrunde, chronologische Sortierung.
  \item Für jede geänderte Liga \code{POST /simulate} mit \num{10000}
        Iterationen, in der 3.~Liga und der 2.~Frauen-Bundesliga zusätzlich
        der Malus-Lauf.
  \item Regionalliga-Kopplung und Aufstiegsspalten berechnen.
  \item Spieldetails je Liga über \code{POST /league-details}.
  \item \code{generate\_static\_site()} schreibt alle Seiten.
\end{enumerate}

\chapter{Der Rust-Simulationsserver}\label{kap:rust}

Die Simulation lief früher in R mit C++-Erweiterungen (Rcpp). Heute
übernimmt sie ein eigener Rust-Server
(\datei{league-simulator-rust/}), den R über HTTP anspricht. Der Server ist
rund hundertmal schneller als der Vorgänger und die \emph{einzige} Stelle,
an der die Modellkonstanten stehen.

\section{Bausteine}

\begin{center}
\small
\begin{tabular}{@{}lll@{}}
\toprule
Crate & Version & Zweck \\
\midrule
\code{axum}       & 0.8.9  & HTTP-Server und Routing \\
\code{tokio}      & 1.53.1 & asynchrone Laufzeit \\
\code{rayon}      & 1.12.0 & parallele Iterationen (Work-Stealing) \\
\code{rand}       & 0.10.3 & Zufallszahlen \\
\code{statrs}     & 0.19.1 & Poisson-Verteilungsfunktion \\
\code{serde}, \code{serde\_json} & 1.0 & JSON \\
\bottomrule
\end{tabular}
\end{center}

Gebaut wird mit \code{rust:1.98-alpine} gegen musl, im Release-Profil mit
\code{opt-level=3}, Link-Time-Optimierung, einer Codegen-Einheit und
\code{panic="abort"}. Der Code gliedert sich in die Module \code{api},
\code{elo}, \code{simulation} (Spiel und Saison), \code{monte\_carlo},
\code{league\_details} und \code{models}.

\section{Endpunkte}

Der Server lauscht auf Port 8080 und nimmt Anfragen bis 2~MB an
(\datei{src/api/mod.rs:15--24}).

\begin{center}
\small
\begin{tabularx}{\textwidth}{@{}lX@{}}
\toprule
Endpunkt & Zweck \\
\midrule
\code{GET /health} & Lebenszeichen für den Docker-Healthcheck: \code{\{"status":"ok", "version":…\}} \\
\code{POST /simulate} & Monte-Carlo-Simulation einer Liga \\
\code{POST /simulate/batch} & mehrere Ligen in einem Aufruf \\
\code{POST /league-details} & deterministischer ELO-Walk und Wahrscheinlichkeiten je Spiel, ohne Zufall \\
\code{POST /match-preview} & ein virtuelles Einzelspiel aus zwei ELO-Werten, für die Aufstiegsspiele \\
\bottomrule
\end{tabularx}
\end{center}

\subsection*{\texttt{POST /simulate}}

\begin{lstlisting}
{
  "schedule":   [[0, 1, 2, 1], [2, 3, null, null], ...],  // Heim, Gast, Tore, Tore
  "elo_values": [2057.26, 1876.14, ...],                    // Start-ELO je Team
  "team_names": ["FCB", "BVB", ...],
  "iterations": 10000,
  "mod_factor": 20,
  "elo_neutral":     [false, true, ...],   // optional, parallel zu schedule
  "adj_points":      [0, -50, ...],        // optional, Malus-Lauf
  "tore_slope": ..., "tore_intercept": ..., // nur Frauen-Ligen
  "group_of_team": [...], "relegation_places": 4   // nur 3. Liga
}
\end{lstlisting}

Mannschaften sind ab 0 nummeriert, fehlende Tore (\code{null}) werden
simuliert. Die Antwort enthält die Wahrscheinlichkeitsmatrix
(Mannschaft~$\times$~Platz, sortiert nach erwartetem Platz), die
Mannschaftsnamen, die Zahl der Simulationen, die Laufzeit und für die
3.~Liga die Zählung \code{relegation\_group\_counts}. Ungültige Anfragen
beantwortet der Server mit Status 400 und einer Klartextmeldung. Mehr als
\num{100000} Iterationen lehnt er ab.

\subsection*{\texttt{POST /league-details}}

Die Anfrage gleicht der von \code{/simulate}, die Mannschaften sind hier
aber ab 1 nummeriert. \code{max\_goals} (Standard 6) bestimmt die Größe der
Ergebnis-Matrix. Die Antwort enthält je Spiel: ob es gespielt ist, die Tore,
die ELO beider Mannschaften vor dem Spiel, die ELO-Änderung, beide
Torerwartungen, $P(\text{Heimsieg})$, $P(\text{Remis})$,
$P(\text{Auswärtssieg})$ und die Ergebnis-Matrix. Dazu kommen die aktuellen
ELO-Werte aller Mannschaften. Aus dieser einen Antwort entstehen Ligatabelle,
Rückblick, Ausblick und ELO-Verlauf.

\section{Wer welche Parameter sendet}

R sendet den Heimvorteil nie, und das Tormodell nur für die Frauen-Ligen.
Fehlt ein Wert, nimmt der Server seinen Standard. Damit gibt es für jede
Konstante genau eine Quelle (ADR~0002). Ein Wächtertest
(\datei{tests/testthat/test-waechter-quelltext.R}) liest den R- und
JavaScript-Quelltext als Text und schlägt fehl, sobald dort eine der
Konstanten auftaucht. R serialisiert JSON mit \code{digits = NA}, damit
\code{tore\_slope} nicht auf vier Nachkommastellen gerundet ankommt.

\section{Leistung}

Gemessen am 7.~September 2026 auf einem Entwicklerrechner, 18 Mannschaften,
voller Spielplan ohne gespielte Partien (der teuerste Fall):

\begin{center}
\small
\begin{tabular}{@{}lr@{}}
\toprule
Messgröße & Wert \\
\midrule
\num{10000} Iterationen & ca. 42~ms \\
\num{50000} Iterationen & ca. 218~ms \\
eine Saison & ca. 4~µs \\
HTTP und JSON auf R-Seite & ca. 20~ms \\
alle zehn Ligen je Zyklus & ca. 0{,}5~s \\
\bottomrule
\end{tabular}
\end{center}

Die Dauer eines Zyklus bestimmen damit die API-Abrufe, nicht die
Simulation.

\chapter{Rendering, Auslieferung, Build und Betrieb}\label{kap:betrieb}

\section{Rendering}

Der Seitengenerator (\datei{RCode/generate_static_site.R}) baut die Seiten
mit dem R-Paket \code{htmltools}. Der Ablauf:

\begin{enumerate}
  \item Nur Ligen, für die alle benötigten Ergebnisse vorliegen, werden
        gerendert. Die übrigen werden übersprungen und im Log genannt: Eine
        Website ohne eine Liga ist besser als gar keine.
  \item Zuerst werden die Verlaufsdaten aller Ligen berechnet. Davon hängen
        der Link auf der Ligaseite und die Ziele der Navigation auf den
        Verlaufsseiten ab.
  \item Danach entstehen die Ligaseiten, die Verlaufsseiten (jede einzeln
        abgesichert), die Aufstiegsseite (nur wenn mindestens eine
        Regionalliga dabei ist) und die Methodik-Seite.
  \item Abschließend steht im Log \enquote{wrote 22 pages}.
\end{enumerate}

\paragraph{Atomares Schreiben.} Jede Seite wird zuerst in eine Datei
\code{<name>.tmp} im selben Verzeichnis geschrieben und dann umbenannt
(\datei{generate_static_site.R:960--980}). Das Umbenennen ist auf demselben
Dateisystem atomar. Der Webserver sieht deshalb immer entweder die alte oder
die vollständige neue Seite, nie eine halb geschriebene.

\paragraph{Assets.} Stilvorlage, Schriften, Favicon und \code{verlauf.js}
werden bei jedem Rendern nach \code{assets/} kopiert.

\section{Auslieferung}

Caddy bindet das Volume nur lesend ein, komprimiert mit zstd und gzip und
setzt Cache-Regeln je Pfad:

\begin{center}
\small
\begin{tabular}{@{}ll@{}}
\toprule
Pfad & \code{Cache-Control} \\
\midrule
Ligaseiten, Verlaufsseiten, Aufstiegsseite & \code{no-cache} \\
\code{methodik.html} & \code{public, max-age=3600} \\
\code{/assets/*} & \code{public, max-age=300} \\
\bottomrule
\end{tabular}
\end{center}

\begin{achtung}[Neue Seiten brauchen eine eigene Caddy-Zeile]
Die Pfadregeln in Caddy gelten exakt. Eine Seite, die im \code{route}-Block
nicht genannt ist, bekommt \emph{gar keinen} Cache-Header. Wer eine Seite
hinzufügt, trägt sie vor dem Deployment in die Caddy-Konfiguration ein. So
geschah es zuletzt am 11.~Oktober 2026 für die zehn Verlaufsseiten.
\end{achtung}

\section{Container}

Das Produktions-Image (\datei{Dockerfile}) wird in drei Stufen gebaut:

\begin{enumerate}
  \item \code{rust:1.98-alpine} kompiliert den Rust-Server.
  \item \code{rocker/r-ver:4.6.1} installiert die R-Pakete mit allen
        Entwicklungsbibliotheken.
  \item Das Laufzeit-Image \code{rocker/r-ver:4.6.1} enthält nur die nötigen
        Bibliotheken, die R-Pakete, das Rust-Programm und \datei{RCode/}.
        Es läuft als Benutzer \code{appuser} (uid 1001), nicht als root.
\end{enumerate}

Die R-Pakete im Produktionspfad sind \code{dplyr}, \code{tidyr},
\code{tibble}, \code{magrittr}, \code{httr}, \code{jsonlite}, \code{digest},
\code{xtable}, \code{here} und \code{htmltools}. Tests und Skripte gehören
nicht zum Image.

\begin{center}
\small
\begin{tabularx}{\textwidth}{@{}lX@{}}
\toprule
Einstellung & Wert \\
\midrule
Healthcheck & \code{curl -f http://localhost:8080/health}, alle 30~s \\
Umgebung & \code{RAPIDAPI\_KEY} (Pflicht), \code{SEASON=2026}, \code{DURATION} leer, \code{TZ=Europe/Berlin}, \code{STATIC\_SITE\_DIR=/app/ShinyApp/public} \\
Volume & \code{fussball-site} $\rightarrow$ \code{/app/ShinyApp/public} \\
Ressourcen & 2~GB Speicher, 4 CPUs, Logs $3 \times 10$~MB \\
Neustart & \code{unless-stopped} \\
\bottomrule
\end{tabularx}
\end{center}

Der Healthcheck prüft nur, ob der Rust-Server lebt. Ein hängender
R-Scheduler macht den Container nicht \enquote{unhealthy}. Nach einem
Deployment muss man deshalb die Logs lesen, nicht nur den Status.

\section{Continuous Integration}

GitHub Actions (\datei{.github/workflows/ci.yml}) prüft jeden Pull Request
und jeden Push auf \code{main}:

\begin{description}[font=\sffamily\bfseries\small]
  \item[rust-quality] \code{cargo fmt --check}, \code{cargo clippy -D warnings},
        \code{cargo test --release} (76 Tests).
  \item[r-lint] \code{lintr} 3.4.0 (fest gepinnt) über \datei{RCode/}, mit
        Zeilenlänge 120. Jeder Befund lässt den Job scheitern.
  \item[image-build-and-test] baut das Produktions-Image, startet darin den
        Rust-Server und führt den vollständigen R-Testsatz aus, einschließlich
        der JavaScript-Tests mit Node 24 und jsdom. Ein Test, der übersprungen
        wird, weil Rust oder Node fehlt, gilt in der CI als Fehler.
  \item[push-image] nur auf \code{main}: Das Image geht mit den Tags
        \code{:latest} und \code{:<kurzer Commit-SHA>} zu Docker Hub
        (\code{chrisschwer/league-simulator}).
  \item[report-failure] legt bei einem Fehlschlag auf \code{main} ein Issue an.
\end{description}

Dazu kommt CodeQL, das wöchentlich und bei jedem Pull Request die
Workflow-Dateien prüft.

\section{Tests}

\begin{itemize}
  \item \textbf{R}: rund 870 \code{testthat}-Blöcke in 61 Dateien. Die
        Dateien folgen einer festen Namensregel
        \code{test-<einheit>[-<thema>].R}, wobei \code{<einheit>} der Name
        einer Datei in \datei{RCode/} ist. Ein Wächtertest erzwingt die Regel.
  \item \textbf{JavaScript}: Sortierung, Tooltips, Veraltet-Hinweis und das
        Verlaufsdiagramm werden in Node mit jsdom getestet, aufgerufen aus den
        R-Tests der erzeugenden Einheit.
  \item \textbf{Rust}: 76 Tests, darunter die Erhaltung der ELO-Summe und
        die Gleichheit von Tormodell in Simulation und Spieldetails.
  \item \textbf{Naht R–Rust}: Integrationstests gegen einen laufenden
        Rust-Server prüfen die JSON-Verträge.
\end{itemize}

\section{Betrieb}

Produktiv läuft der Container auf einem eigenen Server. Dessen
\code{compose.yml} pinnt das Image auf den Commit-SHA eines Merges auf
\code{main}, nie auf \code{latest}, und es gibt keinen automatischen
Updater. Ein Deployment ist im Kern eine einzige geänderte Zeile:

\begin{enumerate}
  \item Umfang prüfen: was sich seit dem laufenden Stand geändert hat, laut
        \code{git log} und nicht laut Pull-Request-Text.
  \item Neue Seiten vorab in die Caddy-Konfiguration eintragen.
  \item Pin auf den neuen SHA setzen, \code{docker compose pull} und
        \code{up -d}.
  \item Verifizieren, und zwar bis über die zweite Runde des Schedulers
        hinaus: kein Neustart-Kreislauf, \enquote{wrote 22 pages}, keine
        Warnungen im Log, alle Seiten mit Status 200 und korrektem
        Cache-Header.
\end{enumerate}

Ein Rollback setzt den Pin auf den vorherigen SHA zurück. Das alte Image
liegt lokal auf dem Server.

\section{Werkzeuge für Betreiber}

\begin{center}
\small
\begin{tabularx}{\textwidth}{@{}lX@{}}
\toprule
Skript & Zweck \\
\midrule
\code{scripts/season\_transition.R} & Saisonwechsel: End-ELO berechnen, neue Stammdatei erzeugen \\
\code{scripts/preview\_site.R} & lokale Vorschau der Website aus einer gespeicherten Ergebnisdatei \\
\code{scripts/verlauf\_fixture.R} & Beispieldaten für die Vorschau der Verlaufsseiten (braucht den Rust-Server) \\
\code{scripts/calibrate\_historical\_elo.R} & einmalige ELO-Eichung neuer Ligen, schreibt nur mit \code{--confirm} \\
\code{scripts/analyze\_league\_empirics.R} & Empirie-Report: Tore je Spiel, Heim/Remis/Auswärts, Remisquoten \\
\bottomrule
\end{tabularx}
\end{center}

Destruktive Werkzeuge laufen standardmäßig als Probelauf und verändern erst
mit \code{--confirm} etwas.


\appendix
\chapter{Formelübersicht}

\begin{center}
\renewcommand{\arraystretch}{1.6}
\begin{tabularx}{\textwidth}{@{}>{\raggedright}p{3.4cm}X@{}}
\toprule
Größe & Formel \\
\midrule
ELO-Erwartung & $p_h = \dfrac{1}{1 + 10^{\clamp_{[-400,400]}(E_a - E_h - H)/400}}$ \\
Score & $S = \tfrac12\bigl(\sgn(g_h - g_a) + 1\bigr)$ \\
Multiplikator & $m = \sqrt{\max(\lvert g_h - g_a\rvert, 1)}$ \\
ELO-Update & $\delta = k\,m\,(S - p_h)$, \quad $E_h' = E_h + \delta$, \quad $E_a' = E_a - \delta$ \\
Torerwartung & $\lambda_{h,a} = \max\bigl(\pm s\,(E_h + H - E_a) + c,\ 0{,}001\bigr)$ \\
Tore & $g \sim \Pois(\lambda)$, unabhängig, per Quantilfunktion \\
Platzwahrscheinlichkeit & $\hat P_{t,p} = C_{t,p}/N$ \\
Platzgruppe & $\sum_{p=l}^{u} \hat P_{t,p}$ \\
RL-Abstieg & $\sum_p \hat P_{t,p}\, w_g(p)$, \quad $w_g(n-d+1) = \sum_{j:\,n_g(j) \ge d} q_g(j)$ \\
RL-Aufstieg (Playoff) & $P_M(X) \sum_Y P_M(Y)\, P(X \succ Y)$ \\
Heatmap-Farbe & $t = p^{0{,}6}$, \quad $\text{RGB} = \operatorname{round}\bigl(255 + t\,(\text{Tinte} - 255)\bigr)$ \\
Standardfehler & $\sqrt{\hat P(1 - \hat P)/N} \le 0{,}005$ \\
\bottomrule
\end{tabularx}
\end{center}

\chapter{Die zehn Ligen}

\begin{center}
\footnotesize
\setlength{\tabcolsep}{4pt}
\begin{tabularx}{\textwidth}{@{}llcll>{\raggedright\arraybackslash}X@{}}
\toprule
Liga & Schlüssel & API-ID & Familie & Teams & Besonderheit \\
\midrule
Bundesliga & \code{bundesliga} & 78 & Herren & 18 & \\
2.~Bundesliga & \code{zweite\_bundesliga} & 79 & Herren & 18 & \\
3.~Liga & \code{dritte\_liga} & 80 & Herren & 20 & Abstieg in fünf Regionen; Malus-Lauf \\
Frauen-Bundesliga & \code{frauen\_bundesliga} & 82 & Frauen & 12--14 & eigenes Tormodell \\
2.~Frauen-Bundesliga & \code{zweite\_frauen\_bl}\textsuperscript{*} & 1034 & Frauen & 14 & drei Absteiger; Malus-Lauf \\
Regionalliga Bayern & \code{rl\_bayern} & 83 & Herren & 16--22 & 2026/27 Aufstiegsspiele; Relegation \\
Regionalliga Nord & \code{rl\_nord} & 84 & Herren & 16--22 & 2026/27 Aufstiegsspiele \\
Regionalliga Nordost & \code{rl\_nordost} & 85 & Herren & 16--22 & 2026/27 Direktaufstieg (Rotation) \\
Regionalliga SüdWest & \code{rl\_suedwest} & 86 & Herren & 16--22 & Direktaufstieg \\
Regionalliga West & \code{rl\_west} & 87 & Herren & 16--22 & Direktaufstieg, Abstieg entkoppelt \\
\bottomrule
\end{tabularx}
\end{center}

\textsuperscript{*}\,Vollständig \code{zweite\_frauen\_bundesliga}. Die Mannschaftszahlen sind Spannen und keine festen Werte. Die Registry
prüft, ob eine Liga in ihrer Spanne liegt, und erlaubt so Ligen, die ihre
Größe ändern.

\chapter{Begriffe}

\begin{description}[font=\sffamily\bfseries\small,style=nextline]
  \item[ELO] Stärkeschätzung einer Mannschaft, Kapitel~\ref{kap:elo}.
  \item[ELO-Walk] Die Kette der ELO-Updates über alle gespielten Spiele einer
        Saison, ausgehend von der Start-ELO.
  \item[Wechselgemeinschaft] Gruppe von Ligen, zwischen denen Mannschaften
        auf- und absteigen und die deshalb ein Tormodell teilen.
  \item[Hauptlauf / Malus-Lauf] Die reguläre Simulation und die zweite
        Simulation, in der nicht aufstiegsberechtigte Zweitvertretungen mit
        $-50$ Punkten starten.
  \item[Abstiegskopplung] Die Abhängigkeit der Regionalliga-Absteiger von den
        Drittliga-Absteigern aus derselben Region.
  \item[Gewertetes Spiel] Am grünen Tisch entschiedenes Spiel (AWD, WO).
        Zählt für die Tabelle, nicht für die ELO.
  \item[Nachholspiel] Ein Spiel, das nach Beginn eines späteren Spieltags
        angepfiffen wird.
  \item[Kürzel] Kurzname einer Mannschaft (\code{ShortText}), nur innerhalb
        einer Liga eindeutig.
  \item[Registry] \datei{RCode/league_registry.R}: alle Ligen und ihre
        Eigenschaften als Daten.
  \item[Statische Seite] Fertige HTML-Datei ohne Server-Logik, Kapitel~\ref{kap:seiten}.
\end{description}

\chapter{Wichtige Quelldateien}

\begin{center}
\small
\begin{tabularx}{\textwidth}{@{}lX@{}}
\toprule
Datei & Inhalt \\
\midrule
\multicolumn{2}{@{}l}{\textit{Rust: Modell}} \\
\datei{src/models/mod.rs} & Modellkonstanten, Datenstrukturen \\
\datei{src/elo/mod.rs} & ELO-Erwartung und Update \\
\datei{src/simulation/match_sim.rs} & Tormodell, Poisson-Ziehung \\
\datei{src/simulation/season.rs} & Saisonsimulation, Tabelle \\
\datei{src/monte_carlo/mod.rs} & Monte Carlo, Zufallszahlen, Parallelität \\
\datei{src/league_details/mod.rs} & deterministischer Walk, 1/X/2, Ergebnis-Matrix \\
\datei{src/api/handlers.rs} & Endpunkte, Prüfung der Anfragen \\
\midrule
\multicolumn{2}{@{}l}{\textit{R: Ablauf und Daten}} \\
\datei{RCode/updateScheduler.R} & Einstieg, Zeitfenster \\
\datei{RCode/update_all_leagues_loop.R} & Hauptschleife, Gating, Haupt- und Malus-Lauf \\
\datei{RCode/retrieveResults.R} & API-Football-Abrufe \\
\datei{RCode/rate_limit_takt.R} & Taktregler \\
\datei{RCode/transform_data.R} & Statusgruppen, Sortierung, Malus \\
\datei{RCode/rust_integration.R} & HTTP-Client für den Rust-Server \\
\datei{RCode/league_details.R} & Spieldetails, Ligatabelle, Fenster \\
\datei{RCode/league_registry.R} & die zehn Ligen als Daten \\
\datei{RCode/rl_abstiegskopplung.R} & Abstiegskopplung \\
\datei{RCode/rl_aufstieg.R}, \datei{aufstiegsspiele.R} & Aufstiegsrotation und -spiele \\
\midrule
\multicolumn{2}{@{}l}{\textit{R und Web: Website}} \\
\datei{RCode/generate_static_site.R} & Seitengenerator \\
\datei{RCode/render_sections.R} & Tabelle, Rückblick, Ausblick, Zonen \\
\datei{RCode/render_helpers.R} & Prozentformat, Heatmap-Farbe, Platzgruppen \\
\datei{RCode/league_views.R} & Ansichten und Spaltengruppen je Liga \\
\datei{RCode/elo_verlauf.R} & Daten der Verlaufsseiten \\
\datei{RCode/site_assets/} & \code{site.css}, \code{verlauf.js}, Schriften, Methodik-Text \\
\midrule
\multicolumn{2}{@{}l}{\textit{Dokumente}} \\
\datei{docs/adr/} & Architekturentscheidungen 0001--0008 \\
\datei{docs/modellannahmen-rl-kopplung.md} & Annahmen der Regionalliga-Kopplung \\
\datei{docs/reports/} & Empirie, Frauen-Tormodell, Heimvorteil \\
\datei{docs/deployment/} & Container, Caddy, Rollback \\
\bottomrule
\end{tabularx}
\end{center}


\end{document}
